\section{AMC 10 10A 2025}
        \begin{enumerate}
        	\item (\textit{AMC 10 10A 2025}) Andy and Betsy both live in Mathville. Andy leaves Mathville on his bicycle at $1:30$, traveling due north at a steady $8$ miles per hour. Betsy leaves on her bicycle from the same point at $2:30$, traveling due east at a steady $12$ miles per hour. At what time will they be exactly the same distance from their common starting point?


 $\textbf{(A)}~3:30\qquad\textbf{(B)}~3:45\qquad\textbf{(C)}~4:00\qquad\textbf{(D)}~4:15\qquad\textbf{(E)}~4:30$

 \textbf{Answer: }E



\textbf{Solution 1}

At $2{:}30$, Andy is $8$ miles ahead. For every hour that they both travel, Betsy gains $4$ miles on Andy. Therefore, it will take her $2$ more hours to be the same distance from the starting point as Andy, which occurs at $\boxed{\textbf{(E) } 4{:}30}$.

 $8h = 12(h-1)$ $\Rightarrow 8h = 12h - 12$ $\Rightarrow 4h = 12$ $\Rightarrow h = 3$

 Since Andy started at $1{:}30$, the catch-up time is $4{:}30$.  Answer: $\textbf{(E)}$.

 Alternatively, from Betsy's perspective:  $8(h+1) = 12h$ $\Rightarrow 8h + 8 = 12h$ $\Rightarrow h = 2$  Same result: $\textbf{(E) } 4{:}30$.

 



\textbf{Solution 2 (Simplest Version)}

Andy goes $8$ mph, Betsy goes $12$ mph.  

 We know that Betsy starts 1 hour after Andy, so she will be 8 miles behind.

 $12x - 8 = 8x$
$4x = 8$ $\,\rightarrow\,$ $x = 2$

 2 hours after Betsy's start time is $\boxed{\text{4:30}}$

 -XaberTooth

 



\textbf{Solution 3}

Try the answer choices:

 At $3{:}30$: Andy $=16$, Betsy $=12$  At $3{:}45$: Andy $=18$, Betsy $=15$  At $4{:}00$: Andy $=20$, Betsy $=18$  At $4{:}15$: Andy $=22$, Betsy $=21$  At $4{:}30$: Andy $=24$, Betsy $=24$You can visualize this using a table:

 \begin{array}{c|c|c}\textbf{Time} & \textbf{Andy - 8 mph} & \textbf{Betsy - 12 mph} \\ \hline3{:}30 & 16 & 12 \\3{:}45 & 18 & 15 \\4{:}00 & 20 & 18 \\4{:}15 & 22 & 21 \\4{:}30 & 24 & 24\end{array}

 Thus the correct answer is $\boxed{\textbf{(E) } 4{:}30}$.

 



\textbf{Solution 4}

Since $\text{lcm}(8,12)=24$, the time for Andy to reach 24 miles is  $24/8 = 3$ hours.

 $1{:}30 + 3{:}00 = \boxed{\textbf{(E) } 4{:}30}$.

 Note: Alternatively we can find that the time for Betsy to reach $24$ miles is $\frac{24}{12} = 2$ hours. So we have

 
 \[2{:}30 + 2{:}00 = \boxed{\textbf{(E) } 4{:}30}.\]




\textbf{Solution 5}

Andy goes $8x$ miles, Betsy goes $12(x-1)$ miles.  Set equal:

 $8x = 12(x-1)$ $\Rightarrow x = 3$

 Thus the answer is $\boxed{\textbf{(E) } 4{:}30}$.

 



\textbf{Solution 6}

Using constant rates:

 
 \[\begin{array}{c|c|c} \textbf{Time} & \textbf{Andy} - 8\text{ mph} & \textbf{Betsy} - 12\text{ mph} \\ \hline \hline 1{:}30 - \text{Starting} & 0 \text{ miles} & 0 \text{ miles} \\ \hline 2{:}30 & 8 \text{ miles} & 0 \text{ miles} \\ \hline 3{:}30 & 16 \text{ miles} & 12 \text{ miles} \\ \hline  4{:}30 & 24 \text{ miles} & 24 \text{ miles} \end{array}\]

 Thus the answer is $\boxed{\textbf{(E) } 4{:}30}$.

 



\textbf{Solution 7 (Calculus Overkill)}

Let  $D_A(t) = 8t$,  $D_B(t) = 12(t-1)$ for $t \ge 1$.

 Difference:  $F(t) = 12(t-1) - 8t = 4t - 12$.

 Set $F(t)=0$:  $4t - 12 = 0$ $t = 3$ hours.

 $1{:}30 + 3 = 4{:}30$.

 -apex304

 



\textbf{Solution 8 (Easy to understand, fast)(Detailed)}

Let \( t \) denote time in hours measured from 1:30 PM:

 \( t := \text{hours after } 1:30 \text{ PM} \).

 Thus \( t = 0 \) at 1:30 PM, \( t = 1 \) at 2:30 PM, and the physically relevant domain for both riders having left is \( t \geq 1 \).

 Velocity/position (vectorized).Choose orthonormal basis \( \{ i, j \} \) with \( i \) east and \( j \) north. Define piecewise position vectors \( r_A(t), r_B(t) \in \mathbb{R}^2 \) (for all \( t \geq 0 \)):

 \( r_A(t) = 8t \, j = (0, 8t) \),
 \[r_B(t) = \begin{cases} (0, 0), & 0 \leq t < 1, \\ 12(t - 1) \, i, & t \geq 1, \end{cases} = (12 \max\{t - 1, 0\}, 0).\]

 Define scalar distance functions (Euclidean norm / 2-norm):

 \( D_A(t) := \| r_A(t) \|_2 = \sqrt{0^2 + (8t)^2} = 8t \),\( D_B(t) := \| r_B(t) \|_2 = 12 \max\{t - 1, 0\} \).

 We seek \( t \) with \( D_A(t) = D_B(t) \) and \( t \geq 1 \). For \( t \geq 1 \) the equality reduces to the linear equation

 \( 8t = 12(t - 1) \).

 Bring like terms together:

 \( 8t = 12t - 12 \implies 12 = 12t - 8t \).

 Compute the difference \( 12t - 8t \) digit-by-digit: \( 12t - 8t = (12 - 8)t = 4t \). So

 \( 12 = 4t \).

 Divide \( 12 \) by \( 4 \) carefully: \( 12 \div 4 = 3 \). Hence

 \( t = 3 \).

 


	\item (\textit{AMC 10 10A 2025}) A box contains $10$ pounds of a nut mix that is $50$ percent peanuts, $20$ percent cashews, and $30$ percent almonds. A second nut mix containing $20$ percent peanuts, $40$ percent cashews, and $40$ percent almonds is added to the box resulting in a new nut mix that is $40$ percent peanuts. How many pounds of cashews are now in the box? 


 $\textbf{(A) } 3.5 \qquad\textbf{(B) } 4 \qquad\textbf{(C) } 4.5 \qquad\textbf{(D) } 5 \qquad\textbf{(E) } 6$

 \textbf{Answer: }B



\textbf{Solution 1}

We are given that $0.2(10) = 2$ pounds of cashews in the first box.

 Denote the pounds of nuts in the second nut mix as $x.$

 \[5 + 0.2x = 0.4(10 + x)\]
 \[5 + 0.2x = 4 + 0.4x\]
 \[5 - 4 = 0.4x - 0.2x\]
 \[0.2x = 1\]
 \[x = 5\]

 Thus, we have $5$ pounds of the second mix. 

 
 \[0.4(5) + 2 = 2 + 2 = \boxed{\text{(B) }4}\]




\textbf{Solution 2}

Let the number of pounds of nuts in the second nut mixture be $x$. Therefore, we get the equation $0.5 \cdot 10 + 0.2 \cdot x = 0.4(x+10)$. Solving it, we then get $x=5$. Therefore, the amount of cashews in the two bags is $(0.2 \cdot 10) + (0.4 \cdot 5) = 4$, so our answer choice is $\boxed{\textbf{(B)} 4}$.

 



\textbf{Solution 3}

The percent of peanuts in the first mix is $10\%$ away from the total percentage of peanuts, and the percent of peanuts in the second mix is $20\%$ away from the total percentage. This means the first mix has twice as many nuts as the second mix, so the second mix has $5$ pounds. $0.20 \cdot 10 + 0.40 \cdot 5 = 4$ pounds of cashews. So our answer is, $\boxed{\textbf{(B)}4}$

 The weight of the second box of nut mix can be calculated by: ($0.2x + 5) / (10 + x) = 0.4$, where the numerator is equal to the weight of the peanuts, the denominator is equal to the total weight of the boxes, and $x$ is equal to the weight of the second box.

 Multiplying both sides by $10 + x$, you get $0.2x + 5 = 4 + 0.4x$

 Isolating the variable: $0.2x = 1$$\Rightarrow x = 5$

 Calculating the weight of the cashews: $5(0.4) + 10(0.2)$$\Rightarrow 2 + 2$$\Rightarrow 4$

 The answer is $\boxed{\textbf{(B) } 4}$.

 



\textbf{Solution 4}

Note that we can set the information given in the problem into a table shown below:

 
 \[\renewcommand{\arraystretch}{1.5} \begin{centering} \begin{array}{| c | c | c |} \hline \text{Peanuts} & \text{Cashews} & \text{Almonds}\\ \hline 5 & 2 & 3\\ \hline \frac{2}{10}x & \frac{4}{10}x & \frac{4}{10}x\\ \hline \end{array} \end{centering}\]

 We are given that the new nut mix will contain $40\%$ peanuts. Hence, $5 + \frac{2}{10}x$ is $40\%$ of the total mix which is $10 + x$. Solving the equation $5 + \frac{2}{10}x = \frac{2}{5} \cdot (10 + x)$ yields $x=5.$ Therefore, the number of cashews in the new mix is equal to $2 + \frac{2}{5} \cdot 5 = \boxed{\textbf{(B)}  4}$.

 



\textbf{Solution 5}

Write composition vectors as column vectors (mass fractions multiplied by mass):

 Initial mass vector (peanuts, cashews, almonds) for 10 lb:
 \[\mathbf{v}_0 = 10 \begin{pmatrix} 0.50 \\ 0.20 \\ 0.30 \end{pmatrix} = \begin{pmatrix} 5 \\ 2 \\ 3 \end{pmatrix}.\]Added mix per unit mass:
 \[\mathbf{u} = \begin{pmatrix} 0.20 \\ 0.40 \\ 0.40 \end{pmatrix}.\]After adding \( x \) lb, total mass vector is
 \[\mathbf{v}(x) = \mathbf{v}_0 + x \mathbf{u} = \begin{pmatrix} 5 \\ 2 \\ 3 \end{pmatrix} + x \begin{pmatrix} 0.2 \\ 0.4 \\ 0.4 \end{pmatrix}.\]Peanut fraction requirement:
 \[\frac{\mathbf{e}_1^\top \mathbf{v}(x)}{\mathbf{1}^\top \mathbf{v}(x)} = 0.4,\]where \( \mathbf{e}_1 = \begin{pmatrix} 1, 0, 0 \end{pmatrix}^\top \) and \( \mathbf{1} = \begin{pmatrix} 1, 1, 1 \end{pmatrix}^\top \). This is exactly the same scalar equation,
 \[\frac{5 + 0.2x}{10 + x} = 0.4,\]leading to \( x = 5 \) and hence
 \[\mathbf{v}(5) = \begin{pmatrix} 5 \\ 2 \\ 3 \end{pmatrix} + 5 \begin{pmatrix} 0.2 \\ 0.4 \\ 0.4 \end{pmatrix} = \begin{pmatrix} 6 \\ 4 \\ 5 \end{pmatrix}.\]Thus cashews \( = 4 \) lb.

 



\textbf{Solution 6: Very easy algebra and straightforward!}

Set up the ratios for the initial and the additional mixtures. The ratio of Peanuts:Cashews:Almonds (denoted as P:C:A) would initially be $5x:2x:3x$ for some factor $x$. Then for the second mixture it would be $2x:4x:4x$. If we add these we get 7x:6x:7x. Now we know that the percentage of peanuts is 40\%. But we also know that the number of pounds of peanuts and almonds is the same. $40\%+40\%=80\%$ and so the number of cashews makes up $20\%$ of the mixture. If you divide the factor $x$ by the whole ratio you get $7:6:7$ . These values sum to $20$ and $20\%$ of $20$ is $4$.

 \textbf{Video Solutions} 


	\item (\textit{AMC 10 10A 2025}) How many isosceles triangles are there with positive area whose side lengths are all positive integers and whose longest side has length $2025$?    


 $\textbf{(A) } 2025 \qquad\textbf{(B) } 2026 \qquad\textbf{(C) } 3012 \qquad\textbf{(D) } 3037 \qquad\textbf{(E) } 4050$

 \textbf{Answer: }D



\textbf{Solution 1: Casework}

You can split the problem into two cases:
Case $1$: The two sides with equal length are both smaller than $2025$, which means that they range from $1013$ to $2024$. There are $1012$ such cases.
Case $2$: There are two sides of length $2025$, so the last side must be in the range $1$ to $2025$. There are $2025$ such cases. Keep in mind, an equilateral triangle also counts as an isosceles triangle, since it has \textbf{at least} 2 sides of equal length.
Therefore, the total number of cases is (equilateral + isosceles) $1012 + 2025 = \boxed{\text{(D) }3037}$
~cw, minor edit by sd

 Side note not by the author: If you're unsure whether the equilateral triangle is a valid case, notice that the other answers are much farther away, as compared to only a one-off error from 3037, which would make 3037 still the best answer.

 



\textbf{Solution 2 (broken down)}

Suppose an isosceles triangle has sides $A, A, B$ where $A$ is the longest side and $A, B$ are positive integers.

 By the triangle inequalities:

 
 \[A + A > B\]

 \[A + B > A \implies B > 0\]

 ---

 Since $A = 2025$ (the longest side):

 
 \[2025 + 2025 > B \implies B < 4050\]

 But also, since $A$ is the longest side, $B < 2025$. Therefore,

 
 \[B = 1, 2, 3, \ldots, 2024,\]

 giving $2024$ triangles.

 ---

 Now consider triangles with side lengths $B, B, 2025$, where $B < 2025$.

 Triangle inequality gives:

 
 \[B + B > 2025 \implies B > 1012.5\]

 Since $B$ is an integer and less than $2025$, the possible values are

 
 \[B = 1013, 1014, 1015, \ldots, 2024,\]

 which gives $1012$ triangles.

 ---

 Finally, don’t forget the case where all sides are $2025$ (equilateral, which also counts as isosceles by definition). So add $1$ more triangle.

 ---

 In total,

 
 \[2024 + 1012 + 1 = \boxed{3037}\]

 triangles.

 



\textbf{Solution 3 (more broken down)}

To solve this question, we first must observe the following restrictions given by the problem -  1.) Be an isosceles triangle  2.) Positive integer side lengths (causing a positive area)  3.) Longest length is $2025$ 

 To satisfy restrictions $1$ and $2$, we can denote the side lengths of a triangle that satisfies these restrictions as $A, A, B,$ where $A$ and $B$ are positive integers. This is possible as the triangle is restricted to be isosceles, meaning at least $2$ equal length sides.

 In any triangle with side lengths $a$, $b$, and $c$, the following inequalities (Triangle Inequality Theorem) must hold:

 
 \[a + b > c\] 
 \[a + c > b\] 
 \[b + c > a\]

 Therefore, our satisfactory triangles must satisfy

 
 \[A + A > B\] 
 \[A + B > A.\]

 where $A > 0, B > 0.$

 Now we can focus on restriction $3.$ Let us consider that $A$ is equal to $2025$ ($A$ is the longest side length).

 Substituting for $A$ we can rewrite our inequalities as follows

 
 \[2025 + 2025 > B\] 
 \[2025 + B > 2025\]

 where $B > 0.$

 Furthermore, we can solve this system of linear inequalities like this

 
 \[4050 > B\] 
 \[B > 0\]

 \[0 < B < 4050\]

 However, for $A$ to be the longest side length, $B$ must be less than $2025$. Therefore, the solutions are

 
 \[\text{integers in } 0 < B < 2025 \implies 1, 2, 3, 4, 5, \ldots, 2024\]

 This creates $2024$ cases.

 ---

 Now let’s consider that $B$ is the longest side length and it is equal to $2025$.

 Substituting for $B$, we can rewrite our inequalities as follows

 
 \[A + A > 2025\] 
 \[A + 2025 > A\]

 where $A > 0.$

 Furthermore, we can solve this system of linear inequalities like this

 
 \[2A > 2025\] 
 \[0 < 2025\] (which can be ignored as it is always true)

 
 \[A > 1012.5\]

 However, for $B$ to be the longest side length, $A$ must be less than $2025$. Therefore, the solutions are

 
 \[\text{integers in } 1012.5 < A < 2025 \implies 1013, 1014, 1015, 1016, \ldots, 2024\]

 This creates $2024 - 1012 = 1012$ cases.

 ---

 But! here is also a case of an isosceles triangle where all side lengths of the triangle are equal in length and congruent - this special isosceles is referred to as an equilateral triangle. We can denote the lengths of such a triangle as $A, A, A$. We need to assign the longest length to $A$, as it is the only place a length can be applied. So, $A = 2025$. We do not need to worry about the Triangle Inequality Theorem as any equilateral shape always satisfies it. For example,

 
 \[A + A > A, \quad A + A > A, \quad A + A > A\]

 where $a,b,c$ of the Triangle Inequality Theorem are all equal to $A$ and $A > 0$

 is always satisfied for any side length, assuming that $A$ is a positive integer.

 This creates $1$ case.

 ---

 Adding together our cases,

 
 \[2024 + 1012 + 1 = \boxed{\text{(D) }3037}.\]

 This is a process called casework.

 



\textbf{Solution 4: Casework}

There are two cases. The first case is where the other two side lengths are the congruent side lengths and less than 2025, and the second one is where 2025 is one of the congruent side lengths.

 Case $1$: We let $a$ be the other side length in the triangle. By the triangle inequality, we have that 
 \[2a>2025 \text{ and } a<2025,\] so we get that $1012.5<a<2025,$ and since $a$ is an integer, $a$ can be $1013,2014,...,2024,$ so there are $1012$ such cases.

 Case $2$: We let $b$ be the other side in the triangle, since 2025 is one of the congruent sides. Since 2025 is the largest side, $a<2025,$ so there are $2024$ solutions in this case.

 However, there is also a case where the triangle is equilateral so we add $1.$

 So, the total number of triangles that are isosceles with largest side $2025$ is $1012+2024+1=\boxed{3037}$



	\item (\textit{AMC 10 10A 2025}) A team of students is going to compete against a team of teachers in a trivia contest. The total number of students and teachers is $15$. Ash, a cousin of one of the students, wants to join the contest. If Ash plays with the students, the average age on that team will increase from $12$ to $14.$ If Ash plays with the teachers, the average age on that team will decrease from $55$ to $52$. How old is Ash?


 $\textbf{(A) } 28 \qquad\textbf{(B) } 29 \qquad\textbf{(C) } 30 \qquad\textbf{(D) } 32 \qquad\textbf{(E) } 33$

 \textbf{Answer: }A



\textbf{Solution 1}

When Ash joins a team, the team's average is pulled towards his age. Let $A$ be Ash's age and $N$ be the number of people on the student team. This means that there are $15-N$ people in the teacher team. Let us write an expression for the change in the average for each team.

 The students originally had an average of $12$, which became $14$ when Ash joined, so there was an increase of $2$. The term $A-12$ represents how much older Ash is compared to the average of the students'. If we divide this by $N+1$, which is the number of people on the student team when Ash joins, we get the average change per team member once Ash is added. Therefore, 
 \[\frac{A-12}{N+1} = 2.\]

 Similarly, for teachers, the average was originally $55$, which decreased by $3$ to become $52$ when Ash joined. Intuitively, $55-A$ represents how much younger Ash is than the average age of the teachers. Dividing this by the expression $(15-N)+1$, which is the new total number of people on the teacher team, represents the average change per team member once Ash joins. We can write the equation

 
 \[\frac{55-A}{16-N} = 3.\]

 To solve the system, multiply equation (1) by $N+1$, and similarly multiply equation (2) by $16-N$. Then add the equations together, canceling $A$, leaving equation $43=50-N$. From this we get $N=7$ and $A= \boxed{28}.$




\textbf{Solution 2}

As shown above, we get the equation 
 \[\frac{55-A}{16-N} = 3.\] Rearranging we get 
 \[\frac{55-A}{3} = 16-N.\] Therefore, since N must be an integer, $55-A$ must be divisible by 3. Plugging in values gets us $A= \boxed{28}.$




\textbf{Solution 3}

Another way is to say that there are $S$ students and $T$ teachers and Ash's age is $A$. We can create the equations\begin{align}&(1) S+T=15 \\&(2) \frac{12S+A}{S+1}=14 \\ &(3) \frac{55T+A}{T+1}=52\end{align}Equation $(2)$ simplifies to $A=2S+14$ and equation $(3)$ simplifies to $A=-3T+52$. Multiply equation $(2)$ by $3$ to get equation \begin{align} (4) 3A=6S+42 \end{align} and equation $(3)$ by $-2$ to get \begin{align}(5) -2A=6T-104 \end{align} Add (4) and (5) to get $A=6(S+T)-62$. After substituting equation (1) and simplifying, you get $A=\boxed{28}$ or answer choice $\boxed{A}$.

 


	\item (\textit{AMC 10 10A 2025}) Consider the sequence of positive integers


 
 \[1,2,1,2,3,2,1,2,3,4,3,2,1,2,3,4,5,4,3,2,1,2,3,4,5,6,5,4,3,2,1,2,\dots\]


 What is the $2025$th term in this sequence?


 $\textbf{(A) } 5 \qquad\textbf{(B) } 15 \qquad\textbf{(C) } 16 \qquad\textbf{(D) } 44 \qquad\textbf{(E) } 45$

 \textbf{Answer: }E



\textbf{Solution 1 (nice equation)}

Group the numbers by their hill pattern: $(12)(1232)(123432)(12345432)...$

 The maximums of each hill occur at terms  $n = 2, 5, 10, 17...$ These terms correspond to maximums of $2, 3, 4,...$ Let $a$ be the maximum at term $N$. Since the sum of the first $x$ odd numbers is $x^2$ we have $1 + (a-1)^2 = N.$ So for example, if a = $4,$ then $N = 10,$ telling us that the peak of the hill with maximum $4$ occurs at the $10$th term.  

 Now, we know $2025 = 45^2,$ so let $a = 46.$ Then $N = 2026,$ so the $2026$th term is $46.$ Then the $2025$th term must be $\boxed{\text{(E) }45}.$




\textbf{Solution 2}

Note that the first $1$ is the first number from the left, the second $1$ is the third number from the left, and the third $1$ is the seventh number from the left. Since the second differences between the indices of the $n^{th}$ $1$’s are equal, there is a quadratic function for the position of the $n^{th}$ $1$. As the second differences are $2$, the leading coefficient of this function is $1$. Let the function be $f(x)=x^2+ax+b$. Then, we have $1+a+b=1$, and $4+2a+b=3$. Solving, we get $a=-1$ and $b=1$. Therefore, the $n$th $1$ from the left is the $(n^2-n+1)$th number from the left. Plugging in $n=45$ gives $f(n)=1981$, so the $1981$st number is a $1$. 

 Also, the sequence shows that the first instance of a positive integer $k$ is always to the left of the $k^{th}$ $1$. This means that following the $1981^{st}$ number, the numbers go to or above $45$. Adding $2025-1981$ to $1$ means that the $2025^{th}$ term of the sequence is $\boxed{\textbf{(E) }45}.$




\textbf{Solution 3 (Quick)}

Note that every $n^2$ th term is $n$. For example, the $3^2$ th term is $3$, the $4^2$ th term is $4$, the $5^2$ th term is $5$, and so on. We also notice that $45^2 = 2025$, so the $2025$th term is $45$, which is answer choice $\boxed{\textbf{(E)}45}.$




\textbf{Solution 4 (Super Simple!!)}

Upon inspecting the problem, we see that the $1^{st}$ term is 1, and so are the $3^{rd}$, $7^{th}$, $13^{th}$, $21^{th}$, $31^{st}$, and so on. We see that the pattern comes from the differences: $+2$, $+4$, $+6$, $+8$, $+10$, and so on. From this we see that a $1$ appears at each position that is a number of the form $n^{2}-n+1$. We know that $2025 = 45 \cdot 45$. Plugging $45$ as $n$ gives $1981$, meaning that the $2025^{th}$ term is $2025-1981+1=45$. Therefore, the solution is $\boxed{\textbf{(E)}45}.$




\textbf{Solution 5}

\(\large \textbf{Takes less than 30 seconds!}\)

 Immediately by looking at the sequence, we observe that the digit in the place of each perfect square is the square root of that number. For example, the first term in the sequence is \(1\) (because \(\sqrt{1}=1\)), the \(4^{th}\) term is \(2\) (because \(\sqrt{4}=2\)), the \(9^{th}\) term is \(3\) (because \(\sqrt{9}=3\)), the \(16^{th}\) term is \(4\) (because \(\sqrt{16}=4\)), the \(25^{th}\) term is \(5\) (because \(\sqrt{25}=5\)), and so on. And since \(2025\) is a perfect square, it makes it easy for us because \(\sqrt{2025}=45\), which makes the answer choice \(\boxed{\textbf{(E) }45}.\)

 To make it easier to understand my way of solving it, here is a table.

 Let \(n =\) the  position of the said number in the sequence

 Let \(\sqrt{n} =\) the number located in that position itself

 \begin{array}{c|c|c}n & \sqrt{n} \\[6pt] \hline1 & 1 \\[6pt]4 & 2 \\[6pt]9 & 3 \\[6pt]16 & 4 \\[6pt]25 & 5 \\[6pt]36 & 6 \\[6pt]\dots & \dots \\[6pt]2025 & \boxed{\textbf{(E) }45}\end{array}

 \(\textbf{Note}\): I am just saying that the square root of the position (if the position is a perfect square) is the number that appears in the said position in the above sequence. I am not saying that the root itself is the first time it appears in the said sequence. To elaborate, the number \(6\) appears on the \(26^{th}\) number, as \(\textit{well}\) as the \(36^{th}\) number as expected (because \(\sqrt{36} = 6\)).

 



\textbf{Solution 6 - Pattern Recognition}

After looking at the problem, it is evident that a pattern exists for the first occurence of a number in the sequence. Knowing this, we can write the places that each number appears at: \begin{align}&1: 1 \\&2: 2 \\&3: 5 \\&4: 10 \\&5: 17 \end{align}If we write down the difference between the each numbers position in sequence, we get $1, 3, 5, \text{and } 7$. We notice how these numbers are similar to how the square numbers increase (square numbers: $1, 4, 9, 16\ldots$). We also notice how the position of the first instance of each number $n$ is just $(n-1)^2+1$. Since we're trying to find the $2025^{th}$ number and $2025 = 45^2$, we can plug in $45$ for $n-1$. As a result, we get that the first instance of the number $46$ appears at $(45)^2+1$ or $2026$. Knowing that the number before the first instance of a number $n$ is $n-1$, we get the answer choice $\boxed{\textbf{(E) }45}.$




\textbf{Solution 7: Patterns}

We list the positions of the $\text{1s}$ in the sequence and get that they appear at $1,3,7,13,...$ We see that, regarding the $n^{th}$ "$1$" in the sequence, the location of the said "$1$" is at the $n(n+1)+1^{th}$ position. So, we can find that the $1$ closest to $2025$ happens where $n=44,$ due to our knowledge that $2025=45^2.$ Since $44(45)+1=1981,$ we can start counting from $1981,$ getting that the $2025^{th}$ term in the sequence is $\boxed{\textbf{(E) 45}}.$




\textbf{Solution 8: Using Arithmetic Sequences}

We can notice that the sequence $1, 2, 1, 2, 3, 2, 1, 2, 3, 4, 3, 2, 1, 2, 3, 4, 5, 4, 3, 2, 1, 2, 3, 4, 5, 6, 5, 4, 3, 2, 1, 2 \dots$ have a special sequence embedded it in. Notice $1, 2$. $1, 2$ have 2 terms, and the next, $1, 2, 3, 2$ have 4 terms. This does not include the $1$ which ends it, which also starts the next one, having a contradiction, so we can construct the following equation, subtracting 1 on the end because the sequence doesn't start with $1, 1, 2, \dots$:

 
 \[2(\frac{n(n+1)}{2})-1=2025\]

 \[n(n+1)-1=2025\]

 Now, we can start to solve the equation:
 \[n(n+1)=2026\]

 \[n^2+n=2026\]Now let us do casework.The closest number to $2026$ that is a square is $2025$, being $45^2$. However, we add the number again, which will result in $2025+45=2070$, over $2026$.We will then do $44^2$, which is $1936$, and adding $44$ to reach the almost end (which is specifically $1, 2, 3, 4, 5, \dots, 44, 45, 44, 43, 42, 41, 40, \dots, 5, 4, 3, 2$) results in $1980$. We add $1$ to make it $1981$, and using the formula $b-a+1$, it is $2025-1981+1$, which equals $\boxed{(E) \:45}$.

 


	\item (\textit{AMC 10 10A 2025}) In an equilateral triangle each interior angle is trisected by a pair of rays. The intersection of the interiors of the middle $20^{\circ}$-angle at each vertex is the interior of a convex hexagon. What is the degree measure of the smallest angle of this hexagon?


 $\textbf{(A) } 80 \qquad\textbf{(B) } 90 \qquad\textbf{(C) } 100 \qquad\textbf{(D) } 110 \qquad\textbf{(E) } 120$

 \textbf{Answer: }C


	\item (\textit{AMC 10 10A 2025}) Suppose $a$ and $b$ are real numbers. When the polynomial $x^3+x^2+ax+b$ is divided by $x-1$, the remainder is $4$. When the polynomial is divided by $x-2$, the remainder is $6$. What is $b-a$?


 $\textbf{(A) } 14 \qquad\textbf{(B) } 15 \qquad\textbf{(C) } 16 \qquad\textbf{(D) } 17 \qquad\textbf{(E) } 18$

 \textbf{Answer: }E



\textbf{Solution 1}

Use synthetic division to find that the remainder of $x^{3}+x^{2}+ax+b$ is $a+b+2$ when divided by $x-1$ and $2a+b+12$ when divided by $x-2$. Now, we solve

 
 \[\begin{cases} a+b+2=4 \\ 2a+b+12 = 6 \\ \end{cases}\]

 This ends up being $a=-8$, $b=10$, so $b-a=10-(-8)=\fbox{\textbf{(E)} 18}$




\textbf{Solution 2}

Via the remainder theorem, we can plug $1$ in for the factor $x-1$ and get $4$, so we have that
 \[1^3+1^2+1a+b=4\]
 \[a+b=2.\]Then, we plug in $2$ and get a remainder of $6$, so we have that
 \[2^3+2^2+2a+b=6\]
 \[8+4+2a+b=6\]
 \[2a+b=-6\]Then, we solve the system of equations.By substitution, we obtain
 \[b=2-a\]
 \[2a+2-a=-6\]
 \[a=-8\]
 \[b=2-(-8)=2+8=10.\]Thus, we have that $b-a=10-(-8)=10+8=\fbox{\textbf{(E)} 18}$




\textbf{Solution 3 (Small)}

Using Remainder theorem, we get that:$4 = (1)^3 + (1)^2 + (1)a + b = 1 + 1 + a + b$ so we get  $2 = a + b$ (i) and from the second statement,$6 = (2)^3 + (2)^2 + (2)a + b = 8 + 4 + 2a + b$ so we get $-6 = 2a + b$ (ii).

 So we Subtract (ii) from (i) to get $a = -8$. Substitute in to get $b = 10$.

 So $b - a = 10 - (-8) = \fbox{\textbf{(E)} 18}$




\textbf{Solution 4 (Clear Concept)}

When the polynomial $f(x)=x^3+x^2+ax+b$ is divided by $x-1$, the remainder is

 \begin{verbatim}
$f(1)= 1+1 + a+ b =4$ ,         $a + b =2$
\end{verbatim}

When the polynomial is divided by $x-2$, the remainder is $6$
\begin{verbatim}
	$f(2)= 8 + 4 + 2a+ b=6$ ,         $2a + b = -6$
\end{verbatim}

Solving ,  $a + b = 2$ and $2a + b = -6$ we get $a = -8$ , $b = 10$ ,hence , $b-a=18$ $\fbox{\textbf{(E)} 18}$



	\item (\textit{AMC 10 10A 2025}) Agnes writes the following four statements on a blank piece of paper.


 $\bullet$ At least one of these statements is true.


 $\bullet$ At least two of these statements are true.


 $\bullet$ At least two of these statements are false.


 $\bullet$ At least one of these statements is false.


 Each statement is either true or false. How many false statements did Agnes write on the paper?


 $\textbf{(A) } 0 \qquad\textbf{(B) } 1 \qquad\textbf{(C) } 2 \qquad\textbf{(D) } 3 \qquad\textbf{(E) } 4$

 \textbf{Answer: }B



\textbf{Solution 1}

We first number all the statements:

 1) At least one of these statements is true.2) At least two of these statements are true.3) At least two of these statements are false.4) At least one of these statements is false.

 We can immediately see that statement 4 must be true, as it would contradict itself if it were false. Similarly, statement 1 must be true, as all the other statements must be false if it were false, which is contradictory because statement 4 is true. Since both 1 and 4 are true, statement 2 has to be true. Therefore, statement 3 is the only false statement, making the answer $\boxed{\text{(B) }1}$. -Rainjs

 



\textbf{Solution 2}

Statements $I,II,$ and $IV$ are true, while statement $III$ is false. Hence, there are $3$ true statements and $\boxed{\text{(B) }1}$ false statement. This result can be checked by examining the statements individually again.

 Statements $I$ and $II$ will be true because there are $3\ge2$ true statements. Statement $IV$ is also true because there is $1\ge1$ false statement. Finally, statement $III$ is false because there are $1\ngeq2$ false statements.

 Note: This solution does not give a path to the solution, but rather only how to check the answer after the answer has been identified.

 



\textbf{Solution 3}

Let's say there are $T$ true statements. We know that $T$ can be any integer from $0$ to $4$. Numbering the $4$ given statements in order as $S1$ to $S4,$ we can represent $S1$ as $T \geq 1$, $S2$ as $T \geq 2$, and since there are $(4-T)$ false statements, $S3$ as $4-T \geq 2 \iff T \leq 2$, and $S4$ as $4-T \geq 1 \iff T \leq 3$.

 Accordingly, we now consider the possible cases for the value of $T$. If $T=0$, then the true statements are $S3$ and $S4$, i.e. there are $2$ true statements in total. But that means $T = 2 \neq 0$, giving a contradiction.

 Similarly, if $T=1$, then the true statements are $S1$, $S3$, and $S4$, i.e. there are $3$ true statements in total. This means $T = 3 \neq 1$, again giving a contradiction.

 If $T=2$, then all $4$ statements are true, yielding the contradiction $T = 4 \neq 2$.

 Finally, however, if $T=3$, then the true statements are $S1$, $S2$, and $S4$. Hence there are again $3$ true statements in total, but now as $T=3$ in the first place, this is actually consistent and not a contradiction.

 For completeness, we also observe that if $T=4$, then the true statements are only $S1$ and $S2$, i.e. $T = 2 \neq 4$, and so this case again leads to a contradiction.

 Thus we conclude that the only possible value is $T = 3$, which gives a consistent solution when $S1$, $S2$, and $S4$ are true and $S3$ is false (as noted above). It follows that the number of false statements is $\boxed{\textbf{(B) }1}$, namely precisely $S3$ (and indeed, $T = 3$ gives $4-T = 4-3 = 1$, as expected).

 



\textbf{Solution 4}

Suppose Statement $I$ is false, then none of the Statements are true, which contradicts the fact that a false Statement $III$ or $IV$ is telling the truth. Therefore, Statement $I$ is true and assume Statement $II$ is false. 

 Statement $II$ thus implies that only Statement $I$ was the truth, and the rest, false. But then, there are 3 false statements but then Statement $III$ and Statement $IV$ are telling the truth. So Statement $II$ is also true. 

 Now, if Statement $III$ is true, then both Statement $III$ and Statement $IV$ is false, contradicting the fact that it is true. Statement $III$ is hence false and Statement $IV$ tells the truth since Statement $III$ lied so indeed, there are at least one lie. There are a total of 3 truths and 1 lie, making the answer $\boxed{\text{(B) }1}$. ~hxve

 



\textbf{Solution 5 (Quick)}

I tried setting that the statements regarding the true ones (I and II) were true and regarding the false ones (III and IV) were false, and it clearly doesn't work.It doesn't work because statement III was true, not false as I set it to be originally.

 What if we make that statement (Statement III) true too? Well, then we have 3 T and 1 F statements, so statement IV is false, statement III is true, and statements I and II are true too.

 So there is $\boxed{\text{(B) }1}$ false statement.

 



\textbf{Solution 6 (Plug in Answer Choices)}

For this question, what I did was to simply test each of the answer choices. 

 Let's start with $\text{(A)}$ $0$ false statements:

 We can check the answer of $0$ false statements against each of the statements.

 Statement $\mathrm{I}$: "At least one of these statements is true."

 This statement is false, which already invalidates the answer of (A) 0 false statements, as we already have a false statement. 

 Let's move onto $\text{(B)}$ $1$ false statement:

 We can check once again this answer against each of the statements.

 Statement $\mathrm{I}$: "At least one of these statements is true."

 This statement is true, as when we have $1$ false statement, we have $3$ true statements, and therefore this is true.

 Statement $\mathrm{II}$: "At least two of these statements are true."

 Once again, this statement is true, as when we have $1$ false statement, we have $3$ true statements, and thus this is also true.

 Statement $\mathrm{III}$: "At least two of these statements are false."

 This statement is false, as we have one false statement, not $2$ and above.

 Statement $\mathrm{IV}$: "At least one of these statements is false."

 This statement is true, as we have one false statement.

 So, after checking the answer choice of $\boxed{\text{(B)} 1}$ false statement across the $4$ statements, we get three true statements and one false statement, making this statement the correct answer. I will not dive into the other $3$ answer choices, as we have already gotten the answer for this question, and if you were doing this way of solving on the test, you would most likely not test the other $3$ answers either for the sake of saving time. 

 



\textbf{Solution 7 (Step by Step, Simple)}

Start by going statement by statement and checking if they can be false. If statement 1 were false statement 4 would be true leading to a contradiction so statement 1 is true. If statement 2 were false 4 would be true and one would be true making statement 2 true. If statement 3 was false and we know statement 1 and statement 2 were true and 4 can be true and it works. Therefore, the correct answer is $\boxed{\text{(B) }1}$.

 --TFEA(would also appreciate if somebody could please latex my proof).--Iceboy1221(i wrote it using latex)

 



\textbf{Solution 8 (for Chinese Committee Paper)}

Due to reasons to prevent cheating, there was a variant for test takers in China, which involves slight iterations in the questions and answer choices. For this particular problem, the question asks for the number of $\mathbf{true}$ statements that Agnes wrote.

 Therefore, utilizing any of the methods in the solutions above, we find that there is one false statement, therefore there are three true statements.

 Therefore, the correct answer choice is $\mathbf{3}$.

 


	\item (\textit{AMC 10 10A 2025}) Let $f(x) = 100x^3 - 300x^2 + 200x$. For how many real numbers $a$ does the graph of $y = f(x - a)$ pass through the point $(1, 25)$?


 $\textbf{(A) }1 \qquad\textbf{(B) }2 \qquad\textbf{(C) }3 \qquad\textbf{(D) }4 \qquad\textbf{(E) }\text{more than }4$

 \textbf{Answer: }C



\textbf{Solution 1}

Substitute $1 - a$ for $x$ and set this expression equal to $25.$ The problem boils down to finding how many real roots  
 \[100(1-a)^3 - 300(1-a)^2 + 200(1-a) = 25\] has. We further simplify this expression and create a function $f(x):$

 \[f(x) = -100a^3 + 100a - 25\] 

 Using Descarte's Rule of Signs we get: 

 Sign changes for $f(x)$ (possible number of positive roots): 2 

 
 \[f(-x) = +100a^3 - 100a - 25\] 

 Sign changes for $f(-x)$ (possible number of negative roots): 1 

 
Possibilities for roots: 

 1) $2$ positive roots, $1$ negative root

 2) $0$ positive roots, $1$ negative root, $2$ imaginary roots

 
So which one is it? We know if the function changes sign between an interval, then a root exists in that interval (Intermediate Value Theorem). From $a = 0$ to $\frac{1}{2},$ the function changes sign because $f(0) = -25$ while $f(\frac{1}{2}) = +\frac{25}{2}$, so a positive root exists. This eliminates the second possibility, implying that there must be $2$ positive and $1$ negative roots. So the answer is $2 + 1 = \boxed{\textbf{(C) } 3}.$ 

 \textbf{Alternative Solution} Since we are trying to find the number of roots of $-100a^3+100a-25$ (which is the same as the number of roots of $4a^3-4a+1$ because we divided by a constant, namely $-25$), we can sketch a graph of this function, which we will call $g(x)$. We find that the function is strictly decreasing when $a<-1$, $g(0)=g(1)=1$, and the function is strictly increasing when $a>1$. We immediately see that since $g(-1)$ is negative, there is a root between $-1$ and $0$ by the Intermediate Value Theorem. From here, we use the first derivative test as follows. Taking the derivative, we have $g'(x)=12a^2-4$. We solve for the roots of this equation, obtaining $\pm \frac{\sqrt{3}}{3}$. We know that at these two values, the function will either reach a local minimum or maximum if the second derivative isn't $0$. Proceeding by the second derivative rule, we have that $g''(x)=24a$ is positive at $a=\frac{\sqrt{3}}{3}$ and negative at $a=-\frac{\sqrt{3}}{3}$, making them local minimums and local maximums, respectively. We don't need to care about the letter case because the local maximum doesn't give any more roots besides the one we've already found. Considering the first case when $a=\frac{\sqrt{3}}{3}$, we must find the sign of $g(a)$. Using some approximations, we have that 
 \[g(\frac{\sqrt{3}}{3})\]
 \[=\frac{4 \sqrt{3}}{9} - \frac{4 \sqrt{3}}{3} + 1\]
 \[\approx \frac{6.9}{9} - \frac{6.9}{3} +1\]
 \[= \frac{23}{30} - \frac{69}{30} + 1\]
 \[= -\frac{16}{30}\]
 \[<0.\]By the Intermediate Value Theorem again, we see that there is one root between $0$ and $\frac{\sqrt{3}}{3}$ and one root between $\frac{\sqrt{3}}{3}$ and $1$. Thus, we have a total of $1+2=3$ roots. Therefore, the answer is $\boxed{\textbf{(C) 3}}$




\textbf{Solution 2}

This problem is essentially asking how many values of $x$ satisfy $f(x)=25,$ since the graph $f(x-a)$ is a shift $a$ units right from the original graph of $f(x).$ Now we just need to determine the amount of real solutions to $100x^3-300x^2+200x=25.$

 Dividing all terms by 25, we get $4x^3-12x^2+8x-1=0.$ Using the rational root theorem and testing, we find no rational roots. Our next method to determine the amount of roots is graphing. Notice how when $x$ $\to$ $+\infty,$ $f(x)$ $\to$ $+\infty.$ When $x$ $\to$ $-\infty,$ $f(x)$ $\to$ $-\infty.$ Substituting $x=0$ gives us $f(x)=-1.$ Substituting $x=1$ gives us $f(x)=-1$ also. This makes us think to see if there was any sudden bump or curve in the graph, so we test the midpoint which is $1/2.$ Substituting $x=1/2,$ we get $1/2-3+4-1=1/2.$ Aha! Connecting the dots, we see there are two solutions from $0\leq x \leq 1.$ Also, $f(x)$ must continue increasing after $x>1,$ which tells us there is a third root when $x>1.$ Now we can conclude that the problem has 3 real roots, meaning that 3 values of a satisfy the problem. The answer is $\boxed {(C) 3}.$




\textbf{Solution 3 (Graphing)}

We can factor $f(x)$ as $100x(x-1)(x-2)$. We just need to check that if the local maximum between $0$ and $1$ is below or above 25. Observe that $f(1/2)>25$, so the graph will cross the line $y=25$ 3 times, $\boxed {C}.$

 \textbf{Somewhat cheese (not recommended)} Instead of plugging in a value between $0$ and $1,$ we can assume that the equation will reach a value above $25$ between those numbers. Thus, the graph crosses $y=25$ 3 times, $\boxed {C}.$

 \textbf{Super Lazy Solution} Instead of using numbers, we could just remember the shape of the graph given, and that it is extremely vertically dilated. The shape of the graph is (for the y-values given) basically $3$ lines, so moving these $3$ lines left and right, we can see that the graph will intersect at 3 points, so $\boxed {C}.$




\textbf{Solution 4 (Discriminant)}

We start by writing $f(x-a)$
 \[f(x-a)=100(x-a)^3-300(x-a)^2+200(x-a)\]Then, we substitute $25$ for $f(x-a)$, and substitute $x=1$
 \[25 = 100(1-a)^3-300(1-a)^2+200(1-a)\]Now, to make this easier, we define $p=1-a$, and find how many real values of $p$ satisfy this
 \[25 = 100p^3-300p^2+200p\]
 \[25 = 100(p^3-3p^2+2p)\]
 \[0.25 = p^3-3p^2+2p\]
 \[0= p^3-3p^2+2p-0.25\]We use the discriminant of a cubic equation, where, assuming the form $ax^3+bx^2+cx+d$, 
 \[\Delta=b^2c^2-4ac^3-4b^3d-27a^2d^2+18abcd\]

 When $\Delta > 0$, there are $3$ distinct real solutions, when $\Delta =0$, there are repeated solutions (which all belong to $\mathbb{R}$), and when $\Delta<0$, there is one real solution and 2 complex-conjugate solutions.

 From the expression we derived, we can see that $a=1$, $b=-3$, $c=2$, and $d=-0.25$

 


 We plug this into the discriminant formula, yielding
 \[\Delta=(-3)^2(2)^2-4(2)^3-4(-3)^3(-0.25)-27(-0.25)^2+18(-3)(2)(-0.25)\]
 \[\Delta=36-32-27-1.6875+27\]
 \[\Delta=4-1.6875\]
 \[\Delta=2.3125\]Therefore, $\Delta>0$, meaning there are three real values of $p$ that satisfy the equation. Since $p=1-a$, we know that there also are 3 real values of $a$ satisfying the equation. Thus,
 \[\Delta > 0 \implies 3 \implies \boxed{C}\]~shockfront99

 



\textbf{Solution 5(Use only if you have like 2 seconds left)}

Ok, so first notice the polynomial is $f(x - a) = 100(x - a)^3 - 300(x - a)^2 + 200(x - a)$. Observe that this is a cubic polynomial and so it has degree $3$ and therefore it has maximum $3$ solutions. Now choose $3$ and move on. :)

 


 \textbf{Note} The above solution doesn't always work as in this problem asked for the number of values for $a$




\textbf{Solution 6 (the most obvious but uses calculus)}

We substitute$\: x-a \:$into the polynomial to obtain

 \begin{align*}f(x-a)&=100(x-a)^3-300(x-a)^2 + 200(x-a) \\&=100(x-a)[(x-a)^2-3(x-a)+2] \\&=100(x-a)(x-a-1)(x-a-2).\end{align*}
We want to find how many values of $a$ exist such that this function goes through $(1,25)$. So we can write
 \[f(1) = 100(1-a)(-1-a)(-a) = 25.\]Expanding gives 
 \[100(1-a)(-1-a)(-a) = 25 \implies 100(-a^2+1)(a) = 25 \implies -a^3+a = \frac{1}{4} \implies a^3-a+\frac{1}{4} = 0.\]So the amount of $a$ such that the function goes through $(1,25)$ is the amount of roots this equation has. From here, we can either try to find all the roots, which I will not try here. Or we can use calculus to find the amount of roots this has. We start by finding the local minimums and maximums, by taking the derivative of this function
 \[\frac{d}{da} (a^3-a+\frac{1}{4}) \implies 3a^2-1.\]Setting the derivative to zero we have
 \[3a^2-1 = 0 \implies 3a^2 = 1 \implies a = \pm \frac{1}{\sqrt{3}}.\]Therefore, the $y$ coordinates of the critical points are
 \[f\left(\frac{1}{\sqrt{3}}\right)=\left(\frac{1}{\sqrt{3}}\right)^{3}-\left(\frac{1}{\sqrt{3}}\right)+\frac{1}{4}=\frac{1}{3\sqrt{3}}-\frac{1}{\sqrt{3}}+\frac{1}{4}=\frac{1-3}{3\sqrt{3}}+\frac{1}{4}=-\frac{2}{3\sqrt{3}}+\frac{1}{4}\]and
 \[f\left(-\frac{1}{\sqrt{3}}\right)=\left(-\frac{1}{\sqrt{3}}\right)^{3}-\left(-\frac{1}{\sqrt{3}}\right)+\frac{1}{4}=-\frac{1}{3\sqrt{3}}+\frac{1}{\sqrt{3}}+\frac{1}{4}=\frac{-1+3}{3\sqrt{3}}+\frac{1}{4}=\frac{2}{3\sqrt{3}}+\frac{1}{4}.\]The first critical point has a negative $y$ because $\frac{2}{3\sqrt{3}}>\frac{1}{4}$ since $\frac{3\sqrt{3}}{2} < 4$ (which are the reciprocals). The second one clearly has a positive $y$ coordinate. The local minimum is negative and the local maximum is positive implies the function crosses the $x$ axis three times implying there are three roots. So there are $\boxed{\textbf{(C) }3}$ values of $a$ such that the function passes through $(1,25)$.

 Remark: If you aren't sure about the positive negative step, then you can verify using quick approximation.

 


	\item (\textit{AMC 10 10A 2025}) A semicircle has diameter $\overline{AB}$ and chord $\overline{CD}$ of length $16$ parallel to $\overline{AB}$. A smaller semicircle
with diameter on $\overline{AB}$ and tangent to $\overline{CD}$ is cut from the larger semicircle, as shown below.


 
                    \begin{figure}[H]
                        \centering
                        \includegraphics[width=0.5\textwidth]{Problem/AMC_10_10A_2025_10_Problem_Figure_1.png}
                    \end{figure}
                    


 What is the area of the resulting figure, shown shaded?


 $\textbf{(A) } 16\pi \qquad\textbf{(B) } 24\pi \qquad\textbf{(C) } 32\pi \qquad\textbf{(D) } 48\pi \qquad\textbf{(E) } 64\pi$

 \textbf{Answer: }C



\textbf{Solution 1 (Somewhat Cheeky)(But Very Effective)}

Notice that the size of the smaller semicircle is not specified, and there is no additional information that hints at any specific size for it. Hence, we can shrink the small semicircle until its area is arbitrarily small and negligible, leaving us with a semicircle with a diameter of $16$. The area of the semicircle is given by $\frac{\pi r^2}{2}$, so we have $r=\frac{16}{2}=8\Rightarrow$$A=\frac{\pi(8)^2}{2}=\boxed{\text{(C) }32\pi}$




\textbf{Solution 2 (Generalized, no assumptions)}

Let the radius of the larger semicircle be $R$ and that of the smaller one be $r.$ We are looking for $\dfrac{1}{2}\pi(R^2-r^2).$ If we connect the center of the large semicircle to one endpoint of the chord and to the center of the chord, we get a right triangle with legs $8$ and $r$ and hypotenuse $R.$ Hence, $R^2=r^2+8^2\implies R^2-r^2=64\implies\dfrac{1}{2}\pi(R^2-r^2)=\boxed{\text{(C) }32\pi}.$




\textbf{Solution 3}

The problem doesn't restrict where the smaller semicircle is along the larger semicircle's diameter. Therefore, we can assume that the two semicircles are concentric. Let the center of both semicircles be $O$, and let $CD$ be tangent to the smaller semicircle at $T$. Let the radius of the smaller semicircle be $x$, and let the radius of the larger semicircle be $r$. If we mirror the diagram over $AB$, we can see that we have two concentric circles. We are trying to find $\pi(\frac{r^2-x^2}{2})$. By Power of a Point on $T$, we can see that
 \[64 = (r + x)(r - x) = r^2 - x^2.\]Thus, $\pi(\frac{r^2 - x^2}{2}) = 32\pi \implies \boxed{\text{(C) }32\pi}.$~vinceS

 



\textbf{Solution 4}

The problem doesn't define where the chord is within the circle. So, let's place the chord such that when lines are drawn from its ends to the center of the larger semicircle, an equilateral triangle with side length $16$ is formed. 

 
                    \begin{figure}[H]
                        \centering
                        \includegraphics[width=0.5\textwidth]{Solution/AMC_10_10A_2025_10_Solution_Figure_1.png}
                    \end{figure}
                    

 In this definition, the radius of the larger semicircle is $16,$ giving it an area of $\frac{\pi\cdot16^2}{2} = 128\pi.$ 

 The height of the equilateral triangle is $8\sqrt{3},$ which is equal to the radius of the smaller semicircle. This gives the smaller semicircle an area of $\frac{\pi\cdot(8\sqrt{3})^2}{2} = 96\pi.$ 

 The total shaded area is the difference of these semicircles, or $128\pi - 96\pi  = \boxed{\text{(C) } 32\pi}.$




\textbf{Solution 5}

Let the radius of the larger semicircle be $R$ and the smaller one $r.$ We are asked to compute $\frac{R^2-r^2}{2} \cdot \pi$. Since a diameter is always perpendicular and bisects a chord, by drawing a diameter and applying Power of a Point Theorem, this yields $(R+r)(R-r)=8^2$ $\implies$ $R^2-r^2=64$. Therefore, the answer is $\frac{64}{2} \pi = \boxed{\text{(C) } 32\pi}.$ ~hxve

 



\textbf{Solution 6}

Since the problem does not restrict where the chord is, WLOG we can simply let the chord be the base of the semicircle. Therefore, the area is simply $\frac{\pi\cdot8^2}{2} = 32\pi.$




\textbf{Solution 7}

We can move this small semicircle to the middle of the big semicircle. Let $r$ be the radius of the small semicircle. By Pythagorean Theorem, the line draw from the midpoint of the diameter of the big semicircle to the big semicircle at the height of the small semicircle (see the diagram) has length $\sqrt{r^2 + 64}$. This is a radius of the big semicircle. So the shaded area would be
 \[\frac{\pi \cdot (\sqrt{r^2+64})^2 - \pi r^2}{2} = \frac{\pi(r^2 + 64 - r^2)}{2} = \frac{64\pi}{2} = \boxed{\text{(C) } 32\pi}.\]Diagram:

 
                    \begin{figure}[H]
                        \centering
                        \includegraphics[width=0.5\textwidth]{Solution/AMC_10_10A_2025_10_Solution_Figure_2.png}
                    \end{figure}
                    




\textbf{Solution 8 (similar to some other stated solutions)}

We see that the chord is very very similar in length to the diameter, so we solve for the area with 16 as the diameter. We find that the area is $64\pi$, and because they want us to subtract the small semicircle and that 16 is actually less than the actual diameter, we can just use half of $64\pi$ as our answer. Thus, the answer is $\boxed{\mathrm{(C)\ }32\pi}$




\textbf{Solution 9}

Let the smaller semicircle $\omega$ be concentric with the larger semicircle $\Omega$. Let $\operatorname{rad}(\omega) = r$ and $\operatorname{rad}(\Omega) = R$. Then, by Pythagoreans, $R^2 = 64+r^2$. By Power of a Point, if $T$ is the tangential point, then $\operatorname{Pow}(T) \implies 64 = (R-r)r$, or $Rr-r^2 = 64$, hence $R^2 = Rr$, so $R(R-r) = 0$. We see $R \neq 0$, so $R = r$. Then, that implies $\omega = \Omega$ hence the chord is actually the diameter. Therefore, the area is $\frac{1}{2} \pi \cdot 64 = \boxed{32\pi}$.

 


	\item (\textit{AMC 10 10A 2025}) The sequence $1,x,y,z$ is arithmetic. The sequence $1,p,q,z$ is geometric. Both sequences are strictly increasing and contain only integers, and $z$ is as small as possible. What is the value of $x+y+z+p+q$?


 $\textbf{(A) } 66 \qquad\textbf{(B) } 91 \qquad\textbf{(C) } 103 \qquad\textbf{(D) } 132 \qquad\textbf{(E) } 149$

 \textbf{Answer: }E



\textbf{Solution 1}

Since the geometric sequence is more restrictive, we can test values for the common ratio until we find one that works. After a few tests, we find that a common ratio of $4$ results in the geometric sequence $1,4,16,64,$ so the arithmetic sequence is $1,22,43,64.$ The answer is $4+16+64+22+43=\boxed{\text{(E) }149}.$

 A more generalized solution is as follows. Let the common difference of the arithmetic sequence be $d$, and the common ratio of the geometric sequence be $r.$ Hence, the two sequences are $1,1+d,1+2d,1+3d$ and $1,r,r^2,r^3.$ Since $z=1+3d=r^3,$ the arithmetic sequence is $1,1+d,1+2d,r^3.$ Since $d=\dfrac{1+3d-1}{3}=\dfrac{r^3-1}{3}$ is a positive integer, we seek the smallest $r\neq1$ such that $r^3-1=(r-1)(r^2+r+1)$ is divisible by $3,$ so the smallest $r$ is $4$. The rest follows like above.

 



\textbf{Solution 2}

Since $1, x, y, z$ is an arithmetic sequence, we have $y = 2x - 1$ and $z = 3x - 2$.Since $1, p, q, z$ is a geometric sequence, we have $q = p^2$ and $z = p^3$.Thus $p^3 = 3x - 2.$

 Because $p^3 \equiv p \pmod{3}$, we get $3x - 2 \equiv p \pmod{3}$, so $p \equiv 1 \pmod{3}$.The smallest integer $p > 1$ satisfying this is $p = 4$.

 Then $64 = 3x - 2$ → $x = 22$, $y = 43$, $z = 64$, $q = 16$.Therefore, $x + y + z + p + q = 22 + 43 + 64 + 4 + 16 = \boxed{\text{(E) }149}$.

 



\textbf{Solution 3 (Satisfying and Simple)}

So we have that the arithmetic sequence is $1, x, 2x-1, 3x-2$ sincecommon difference is $x - 1$.

 Similarly, geometric sequence is $1, p, p^2, p^3$

 So $p^3 = 3x-2$. We test values above 1 (increasing sequence)

 We see that $p = 4, x = 22$ so our answer is $22 + 43 + 64 + 4 + 16 = \boxed{149}$




\textbf{Solution 4 (Easy to identify)}

Since only integer ratios work for the geometric series, test the ratio of $2$, $3$, and $4$, by having $z$ be $2$, $3$, and $4$ cubed, and finding that $\frac{64-1}{3}=21$ (since that is the common ratio of the arithmetic sequence to get $64$ from $1$ in $4$ terms), then $22+43+4+16+64=149$, which is E.

 Note: We divide $z$ by $3$ because to get from $p$ to $z$ you have to add $3d$ which is the difference of the arithmetic sum.

 -Amon26(note)

 



\textbf{Solution 5 (Ratio)}

We have $1, x, y, z$ as an arithmetic sequence. If we let $d$ be the rate of change, $z - y = d$, $y - x = d$ and $x - 1 = d$, now you could follow with solution 3, or, you could just have the geometric sequence ratio be $r$, so we have $\frac{z}{q} = r$, $\frac{q}{p} = r$, $\frac{p}{1} = r$.

 Wait, we see $\frac{p}{1} = r$, so $p = r$, then $q = p^2$, then $p^3 = z$, and since these have to be integers, we have to get a value $p$ that works as an integer for the arithmetic sequence.

 Because we add our common difference $d$ 3 times, we see that $(p^3-1)/3=d$ where d is an integer (this is because if we try 3 we get $z = p$ which violates $p^3 = z$, where $p = 3$, however, the starting is one, so we find a value that when subtracted by 1, gets us a multiple of 3 so the arithmetic sequence works) and we find a value $p = 4$.

 The arithmetic sequence is $1, 22, 43, 64$, where $21 = d$ and it follows as $p^3 - 1 = 63 = 3d$, and adding it up we get 129 for the arithmetic and since 64 is already $z$ so we only need to add $p$ and $q$ or $p^2$, we would get 20 and $20 + 129 = \boxed{\textbf{(E) } 149}$.

 - Randomuser456(Note there are definitely other values where $(p^3-1)/3=d$ is true, however the problem says for the lowest value of z) and this is also my first solution, please change things if necessary:)

 \begin{verbatim}
\textasciicircum{}\_\textasciicircum{} /|\  👍  | / \
\end{verbatim}




	\item (\textit{AMC 10 10A 2025}) Carlos uses a $4$-digit passcode to unlock his computer. In his passcode, exactly one digit is even, exactly one (possibly different) digit is prime, and no digit is $0$. How many $4$-digit passcodes satisfy these conditions?


 $\textbf{(A) } 176 \qquad\textbf{(B) } 192 \qquad\textbf{(C) } 432 \qquad\textbf{(D) } 464 \qquad\textbf{(E) } 608$

 \textbf{Answer: }D



\textbf{Solution 1: CRAZY and Super Fast Casework}

The only two digits that are neither prime nor even are $1$ and $9$. We split this problem into cases based on the number of $2$s. This is because $2$ is both a prime number and an even number. 

 Case 1: In this case, there are no $2$'s. There are $4$ choices for where the even digit goes, and $3$ choices for what the even digit is. There are then $3$ choices for where the prime digit goes, and $3$ choices for what the prime number digit is. The last two spots have $2$ choices each, either $1$ or $9$. This gives a total of $4\cdot 3^3 \cdot 2^2 = 432$ options for this case. 

 Case 2:  In this case, there is one $2$, with it being the only even prime number. There are $4$ choices for where $2$ goes, and $2$ choices for the other three digits each. This case gives a total of $2^3\cdot 4 = 32$ options.  

 Hence, the answer is $432 + 32 = \boxed{\textbf{(D) }464}$ ~Tacos\_are\_yummy\_1

 



\textbf{Solution 2: Cheese but fast}

Let us count the cases where there is a $2$ anywhere in the lock. There are $4$ places to arrange the $2$s, and the remaining digits can only be $1$s and $9$s (since they neither can be even nor prime). Thus, there will be $4\cdot2^3=32$ choices for there to exist one $2$. We note that answer choice $\text{(D)}$ is $32$ greater than another answer choice (specifically, $\text{(C)}$) and is the only answer choice to have this property. Assuming that the pitfall of forgetting the case with a $2$ would be a separate answer choice on its own, we can justify that the answer is $\boxed{\text{(D) }464}$



	\item (\textit{AMC 10 10A 2025}) In the figure below, the outside square contains infinitely many squares, each of them with the same center and sides parallel to the outside square. The ratio of the side length of a square to the side length of the next inner square is $k$, where $0 < k < 1.$ The spaces between squares are alternately shaded as shown in the figure (which is not necessarily drawn to scale). 


 The area of the shaded portion of the figure is $64\%$ of the area of the original square. What is $k$?


 
                    \begin{figure}[H]
                        \centering
                        \includegraphics[width=0.5\textwidth]{Problem/AMC_10_10A_2025_13_Problem_Figure_1.png}
                    \end{figure}
                    
$\textbf{(A) } \frac{3}{5} \qquad\textbf{(B) } \frac{16}{25} \qquad\textbf{(C) } \frac{2}{3} \qquad\textbf{(D) } \frac{3}{4} \qquad\textbf{(E) } \frac{4}{5}$

 \textbf{Answer: }D



\textbf{Solution 1}

Let the side length of the largest square be $a,$ so it has area $a^2.$ Hence, the second-largest square has area $a^2k^2,$ the third-largest has $a^2k^4,$ and so on. 

 It follows that the total shaded area is 
 \[a^2-a^2k^2+a^2k^4-a^2k^6+...=a^2(1-k^2+k^4-k^6+...)=a^2\dfrac{1}{1+k^2}.\]

 The ratio of the area of the shaded region to that of the original square is then 
 \[\dfrac{a^2\frac{1}{1+k^2}}{a^2}=\dfrac{1}{1+k^2}=\dfrac{64}{100}\]
 \[\implies 64+64k^2=100\implies k^2=\dfrac{36}{64}\implies k=\boxed{\text{(D) }\dfrac{3}{4}}.\]

 \textbf{Remark} We can just let $a=1$ because the question deals with ratios, meaning that there wouldn't be a loss of generality if we let the side length equal some value, getting the same answer $\boxed{D}$. This was done in solution 3.

 



\textbf{Solution 2}

Let the side length of the first square be $a$ and the second square be $b$. The area of the original square is $a^2.$ The area of the outermost shaded region is $a^2 - b^2.$ Let $\frac{a}{b} = r.$ We have $b = \frac{a}{r}.$ So the outermost shaded region area becomes: 

 
 \[a^2 - b^2 = a^2 - \frac{a^2}{r^2}\]

 Now let the next square's side lengths be $c$ and $d.$ Similarly, $c = \frac{b}{r} = \frac{a}{r^2}$ and $d = \frac{c}{r} = \frac{a}{r^3}$and the area of the next shaded region becomes: 
 \[c^2 - d^2 = \frac{a^2}{r^4} - \frac{a^2}{r^6}\] 

 Notice the pattern of adding $4$ to the exponent. If this sequence continues infinitely, we ultimately get: 

 
 \[(a^2 + \frac{a^2}{r^4} + \frac{a^2}{r^8} + ...) - (\frac{a^2}{r^2} + \frac{a^2}{r^6} + ...)\] 

 Which can be simplified using the infinite geometric sequence formula. The problem also tells us that all of this equals $64\% a^2,$ or $\frac{16}{25}a^2:$

 \[\frac{a^2}{1 - \frac{1}{r^4}} - \frac{\frac{a^2}{r^2}}{1 - \frac{1}{r^4}} = \frac{16}{25}a^2\]

 \[\frac{1}{1 - \frac{1}{r^4}} - \frac{\frac{1}{r^2}}{1 - \frac{1}{r^4}} = \frac{16}{25}\]

 \[r^2 = \frac{16}{9} \implies r = \frac{4}{3}\] Since the problem tells us the ratio must be less than one, we take the reciprocal to finally get  
 \[k = \frac{1}{r} = \boxed{\text{(D) }\dfrac{3}{4}}.\]




\textbf{Solution 3 (Faster)}

Let the outside square be of side length $1$, and let $p = k^2$. Then the area of the square is $1-p+p^{2}-p^{3}+p^{4}-p^{5} \dots$ or $\left(1-p\right)\left(1+p^{2}+p^{4}+p^{6}+\dots\right)$. Then, using geometric series, we get $\left(1-p\right)\left(\frac{1}{1-p^{2}}\right) = \frac{1}{1+p}$. This is equal to $\frac{64}{100}$. Therefore, $1+p=\frac{25}{16}$, so $p = 9/16$ and $k = \frac{3}{4}$. Therefore, the answer is $\boxed{(D)\frac{3}{4}}.$




\textbf{Solution 4 (The fastest, no sequences/series)}

Since this is an infinite fixed pattern, we can actually take a huge shortcut. Suppose we take off the outer "layer" and then swap the colors. We have essentially created a smaller copy of the original diagram, specifically with a $k:1$ ratio of side lengths. Assume the original diagram had 100 units of area. Then, based on the givens, the original diagram had 64 units of shaded area and 36 units of unshaded area, meaning our new smaller diagram has 36 units of shaded area. Since the ratio of corresponding side lengths between our new diagram and the original diagram is $k$, the ratio of corresponding areas is $k^2$. Thus, we get $k^2 = \frac{36}{64}$, so $k = \boxed{\text{(D) }\dfrac{3}{4}}$.

 Note: the explanation seems complicated to explain, but the thought process while doing it is extremely fast, much faster than any of the methods involving a series.

 


 



\textbf{Solution 5 (based very closely on solution 4)}

Similarly to solution 4, we can imagine "peeling" off the outer layer to get an inner figure that is similar to the outer figure. Letting the large square have side 1, we know the outer shaded layer has area $1-k^2$. The shaded part of the inner figure (with the outer shaded part removed now must make up $\frac{36}{100} = \frac{9}{25}$ of the inner area of $k^2$. Together, the inner shaded part and the outer shaded part must add up to $\frac{64}{100}$, as they make up 64\% of the square. Thus, we get the equation$1-k^2 + \frac{9k}{25} = 16/25$- Solving this gives k = 3/4~Neber100

 



\textbf{Solution 6 (The longest way)}

Notes that: The sum $S$ can be written as:
 \begin{align*}     S &= 1^2 - k^2 + (k^2)^2 - (k^3)^2 + (k^4)^2 - (k^5)^2 + (k^6)^2 - \dots \\     &= k^0 - k^2 + k^4 - k^6 + k^8 - k^{10} + \dots \\     &= (k^0 + k^4 + k^8 + k^{12} + \dots) - (k^2 + k^6 + k^{10} + k^{14} + \dots) \\ \end{align*}
So we have: 
 \begin{align*}     S &= \lim_{n\rightarrow\infty}\left[1\times\dfrac{1-(k^4)^n}{1-k^4}-k^2\times\dfrac{1-(k^4)^n}{1-k^4}\right]\\     &= \dfrac{1}{1-k^4}-\dfrac{k^2}{1-k^4} \\     &= \dfrac{1-k^2}{1-k^4} \\     &= \dfrac{1-k^2}{(1+k^2)(1-k^2)} \\     &= \dfrac{1}{1+k^2} \end{align*}
We are given the condition that:
 \[\because \dfrac{1}{1+k^2}=64\%=\dfrac{64}{100}=\dfrac{16}{25}\]We can now solve for $k$:
 \begin{align*}   \therefore \dfrac{1}{1+k^2} &= \dfrac{16}{25} \\   1+k^2 &= \dfrac{25}{16} \\   k^2 &= \dfrac{25}{16} - 1 \\   k^2 &= \dfrac{25}{16} - \dfrac{16}{16} \\   k^2 &= \dfrac{9}{16} \\   k &= \pm\sqrt{\dfrac{9}{16}} \\   k &= \pm\dfrac{3}{4} \end{align*}
Therefore, ignore the negative value (Because \(0<k<1\)), so the final answer is $\color{red}\boxed{\color{black}\dfrac{3}{4}}\color{black}$. We are done ~~ : )

 



\textbf{Solution 7}

Assume that the side length for the largest square in the diagram is $1.$ Since the ratio of the next inner square to the largest square is $k,$ we get that the side length for the second largest square is $k,$ $k^2$ for the third largest square, and $k^3$ for the fourth largest square. So, we can write that 
 \[(1-k^2)+(k^4-k^6)+...=\frac{64}{100}\] by using the diagram since the area of a square with side length $s$ is $s^2.$ This is a geometric sequence with first term $1-k^2$ and ratio $k^4.$ So, we can sum the geometric series for 
 \[\frac{1-k^2}{1-k^4}=\frac{64}{100}.\] Cross multiplying, we have that 
 \[100-100k^2=64-64k^4.\] Moving the terms to one side, we have $64k^4-100k^2+36=0,$ and we can factor as 
 \[(64k^2-36)(k^2-1)=0.\] Since $k<1,$ $k^2-1=0$ is invalid. So, we have 
 \[64k^2-36=0,\] and simplifying we get $k^2=\frac{36}{64}.$ Since $k>0,$ 
 \[k=\frac{6}{8}=\boxed{\frac{3}{4}}.\]




\textbf{Solution 8 (Cheese )}

Measure the side length of the first square to be 2 inches. Measure the side length of the second square to be 1.5 inches. 
 \[\frac{1.5}{2} = \frac{3}{4} = \boxed{\text{(D) } \frac{3}{4}}\]




\textbf{Solution 9 (Fastest and shortest, similar to solution 4, no series)}

Since the figure is self-similar, the nonshaded area is $100\%-64\% = 36\%$ of the largest square and $64\%$ of the square with side length $k$, which has area $k^2$. Thus $36\% = 64\%k^2$, $k^2 = \frac{36}{64}$, so $k = \frac{6}{8} = \boxed{\text{(D) }\dfrac{3}{4}}$.

 


 



\textbf{Solution 10 (Series Subtraction)}

Assume initial length $1$ since the problem just uses ratios. We see that each side length shrinks by factor $k$, then they are squared to get the areas, which really just multiplies their exponents by $2$. We add our shaded areas and subtract our un-shaded areas, yielding the alternating series 
 \[1 - k^{2} + k^{4} - k^{6}...\]We express this as two infinite series, one being subtracted from the other
 \[\sum_{n=0}^{\infty} k^{4n} -\sum_{n=0}^{\infty} k^{4n+2}\]Rewrite using the infinite geometric series formula and simplify as far as possible

 
 \[\frac{1}{1-k^{4}} - \frac{k^{2}}{1-k^{4}}\]
 \[\frac{1-k^{2}}{1-k^{4}}\]
 \[\frac{1-k^{2}}{(1+k^{2})(1-k^{2})}\]
 \[\frac{1}{1+k^{2}}\]We know this is equal to $\frac{64}{100} = \frac{16}{25}$
 \[\frac{1}{1+k^{2}} = \frac{16}{25}\]
 \[16k^{2}+16 =25\]
 \[16k^{2} = 9\]
 \[k^{2} = \frac{9}{16}\]We must only consider the positive root
 \[k=\boxed{\text{(D) }\dfrac{3}{4}}.\]~shockfront99

 


	\item (\textit{AMC 10 10A 2025}) Six chairs are arranged around a round table. Two students and two teachers randomly select four of the chairs to sit in. What is the probability that the two students will sit in two adjacent chairs and the two teachers will also sit in two adjacent chairs?


 $\textbf{(A) } \frac{1}{6} \qquad\textbf{(B) } \frac{1}{5} \qquad\textbf{(C) } \frac{2}{9} \qquad\textbf{(D) } \frac{3}{13} \qquad\textbf{(E) } \frac{1}{4}$

 \textbf{Answer: }B



\textbf{Solution 1 (20 seconds)}

Pair two students together and put them adjacent on any two seats. There are 6 ways to do this. Considering one of these cases (they are all the same), there are 4 seats left, in which we wish to arrange the teachers together. So pair the teachers together and put them adjacent on any two seats not already occupied by two of the students. There are 3 ways to do this. For all 6 cases, there are \( 6 \times 3 = 18 \) favorable outcomes.

 The number of ways to arrange the 2 students and 2 teachers is $\binom{6}{2} \times \binom{4}{2} = 90$.

 Our probability is \( \frac{18}{90} = \) $\boxed{\textbf{(B) } \frac 15}$




\textbf{Solution 2}

We first count the number of ways to place $4$ distinct people into $6$ distinct chairs: $6\cdot5\cdot4\cdot3 = 360$.

 We now count how many favorable assignments there are. There are $6$ ways to choose an adjacent pair of chairs for the two students. We can arrange the two students in those two chairs in $2$ ways.

 After those two adjacent chairs are taken, there are $4$ chairs left, with $3$ adjacent pairs among them. We choose one of these pairs, arranging the teachers for $3 \cdot 2 = 6$ ways.

 There are $6\cdot 2 \cdot 6 = 72$ favorable arrangements.

 The probability is therefore $\frac{72}{360} = \boxed{\frac{1}{5}}$




\textbf{Solution 3}

We split the problem into cases of where the first teacher and first student sits. The first student sits in a seat with $1$ probability

 


 \textbf{Case $1$:} The first teacher sits next to the first student. There's are $2$ ways to do this so this will happen with $\frac{2}{5}$ chance. Now, there is one valid seat for the second student to sit in with probability $\frac{1}{4}$ and and one valid valid seat for the second teacher to sit in with probability $\frac{1}{3}$. Total probability of this case is $\frac{2}{5}\cdot \frac{1}{4}\cdot \frac{1}{3} = \frac{1}{30}$

 


 \textbf{Case $2$:} The first teacher sits on the opposite end of the circle with $1/5$ chance. This means the second student and the second teacher each have $2$ valid spots to sit in, with probability $\frac{2}{4}$ and $\frac{2}{3}$. Total probability is $\frac{1}{5}\cdot \frac{1}{2}\cdot \frac{2}{3} = \frac{1}{15}$.

 


 \textbf{Case $3$:} The first teacher sits one chair away from the first student with $\frac{2}{5}$ probability because there are $2$ seats that are one chair away. 

 Case $3.1$:

 The second student sits in between the first teacher and first student with $\frac{1}{4}$ chance. The second teacher only has one valid seat, with $\frac{1}{3}$ chance to sit in it. Total probability is $\frac{1}{4}\cdot \frac{1}{3} = \frac{1}{12}$.

 Case $3.2$:

 The second student doesn't sit in between the teacher and the student with a $\frac{1}{4}$ chance. The second teacher can sit in $2$ valid seats with a $\frac{2}{3}$ chance. Total probability is $\frac{1}{4}\cdot \frac{2}{3} = \frac{1}{6}$.

 Total probability for case $3$ is $\frac{2}{5}\cdot (\frac{1}{12} + \frac{1}{6}) = \frac{1}{10}$.

 
Adding up all the cases, you get $\frac{1}{30} + \frac{1}{15} + \frac{1}{10} = \boxed{\frac{1}{5}}$.

 



\textbf{Solution 4}

We consider how many ways there are to place the students and teachers at random at the table. This is a word arrangement in a circle, and so there are $\frac{(6-1)!}{2! \cdot 2!\cdot 2!}=15$ different ways of doing so, assuming that teachers, students and empty spaces are indistinguishable between themselves. 

 Now, we wish to find how many favorable arrangments there are, so suppose we fix two adjacent teachers. Then, there are three "spots" (blocks of two seats) two adjacent students could occupy: two on the side of either teacher, and one such that no student is sitting adjacent to a teacher.

 Thus, the probability is $\frac{3}{15}=\boxed{\frac{1}{5}}$.

 



\textbf{Solution 5 (Step-by-step, needs verification)}

Because this is in a circle, rotating around doesn't increase or decrease the probability (so pin one seat and the others will not rotate).

 Hence, label the spots 1-6 in clockwise. Let student 1 sit in spot 1. Student 2 can therefore sit in either spot 2 or 6, which the probability is 2/5. For explanation purposes, put student 2 in position 2. 

 Here then requires a case breakdown. 

 Case 1: Teacher $1$ sits in spot $3$, hence teacher $2$ needs to sit in spot $4$ to satisfy. This case totals up with $\frac{2}{5}\cdot \frac{1}{4} \cdot \frac{1}{3} = \frac{1}{30}$.

 Case 2: Teacher $1$ sits in spot $4$, hence teacher $2$ needs to sit in spot $3$ or $5$ to satisfy. This case totals up with $\frac{2}{5} \cdot \frac{1}{4} \cdot \frac{2}{3} = \frac{1}{15}$.

 Case 3: Teacher $1$ sits in spot $5$, hence teacher $2$ needs to sit in spot $4$ or $6$ to satisfy. This case totals up with $\frac{2}{5} \cdot \frac{1}{4} \cdot \frac{2}{3} = \frac{1}{15}$.

 Case 4: Teacher $1$ sits in spot $6$, hence teacher $2$ needs to sit in spot $5$ to satisfy. This case totals up with $\frac{2}{5} \cdot \frac{1}{4} \cdot \frac{1}{3} = \frac{1}{30}$.

 By adding the case probabilities, we have $(\frac{1}{30})\cdot 2 + (\frac{1}{15}) \cdot 2 = \boxed{\frac{1}{5}}$.

 Note: Feel free to delete if there is a logic error.~Mitsuihisashi14

 



\textbf{Solution 6 (Quickest)}

For the two teachers' seatings, there are $6$ ways that they are adjacent out of a combination of $n = \binom{6}{2}$ ways. This probability is $\frac{6}{15} = \frac{2}{5}$.For the two students' seatings, there are $3$ ways to choose the two seats out of $n = \binom{4}{2}$ ways. This probability is $\frac{3}{6} = \frac{1}{2}$ . Answer is $\frac{2}{5} \cdot \frac{1}{2} = \boxed{\frac{1}{5}}$.

 



\textbf{Solution 7 (Probability = Desired Arrangements/Total arangements)}

Let the seats be numbered $1,2,3,4,5,6$ in clockwise order. Then assume placing a group of $2$ (either the two students of the two teachers) in seats $1$ and $2$, this group of two can be permuted (or in this case just swpped) in $2!=2$ ways. Now the second group may sit in the seats $3,4$ or $4,5$ or $5,6$ giving a factor of $3$. This group can also be permuted in $2!=2$ ways. This gives 12 ways for the 4 people to be seated under the condition. Without the condition there are $\dfrac{(6-1)!}{2}=60$ ways for the 4 people to be seated (the division by 2! comes from the two empty seats being considered the same). Finally, the desired probability is:
 \[\frac{12}{60}=\boxed{\frac{1}{5}}\]




\textbf{Solution 8}

For the two students, you can seat the first one in any seat, followed by the second one with a $\frac{2}{5}$ possibility of satisfying the condition. After this, we have two cases.

 Case 1: The first teacher sits next to the two students. This can happen $\frac{1}{2}$ of the time and the other teacher has a $\frac{1}{3}$ chance of sitting next to the first one. This gives us a $\frac{1}{2} \cdot \frac{1}{3} = \frac{1}{6}$ chance for this case.

 Case 2: The first teacher sits away from the two students. This can also happen $\frac{1}{2}$ of the time and the other teacher has a $\frac{2}{3}$ chance of sitting next to the first one. This gives us a $\frac{1}{2} \cdot \frac{2}{3} = \frac{1}{3}$ chance for the second case.

 Now, we can find the desired probability as: $(\frac{1}{3} + \frac{1}{6}) \cdot \frac{2}{5} = \frac{1}{2} \cdot \frac{2}{5} = \boxed{\frac{1}{5}}$



	\item (\textit{AMC 10 10A 2025}) In the figure below, $ABEF$ is a rectangle, $\overline{AD}\perp\overline{DE}$, $AF=7$, $AB=1$, and $AD=5$.
What is the area of $\triangle ABC$?


 
                    \begin{figure}[H]
                        \centering
                        \includegraphics[width=0.5\textwidth]{Problem/AMC_10_10A_2025_15_Problem_Figure_1.png}
                    \end{figure}
                    
$\textbf{(A) } \frac{3}{8} \qquad\textbf{(B) } \frac{4}{9} \qquad\textbf{(C) } \frac{1}{8}\sqrt{13} \qquad\textbf{(D) } \frac{7}{15} \qquad\textbf{(E) } \frac{1}{8}\sqrt{15}$

 \textbf{Answer: }A



\textbf{Solution 1}

Because $ABEF$ is a rectangle, $\angle ABC=90°$. We are given that $\angle ADE=90°$, and since $\angle ECD=\angle ACB$ by vertical angles, $\triangle ECD \sim \triangle ACB$. Let $AC=x$. By the Pythagorean Theorem, $CB=\sqrt{x^2-1}$. Since $AF=BE=7$, $EC=7-\sqrt{x^2-1}$. Because $AC=x$ and $AD=5$, $CD=5-x$. By similar triangles, 
 \[\frac{7-\sqrt{x^2-1}}{x}=\frac{5-x}{\sqrt{x^2-1}}\]. Cross-multiplying, we get that 
 \[7\sqrt{x^2-1}-x^2+1=5x-x^2\], so 
 \[7\sqrt{x^2-1}=5x-1\]. We square both sides, and this is simply a quadratic in $x$: 
 \[24x^2+10x-50=0\], which has a positive root $x=\frac{5}{4}$. Since $AB=1$, we can plug this into the Pythagorean Theorem, with $\frac{5}{4}$ being the hypotenuse, to get $BC=\frac{3}{4}$, and ${1}\cdot \frac{\frac{3}{4}}{2}$ to equal $[ABC]= \boxed{\textbf{(A)} \frac{3}{8}}$

 Solution by HumblePotato, written by lhfriend, 

 Minor edit by SixthGradeBookWorm927

 Minor edit by aldzandrtc

 Minor edit by rcll (I fixed this, as aldzandrtc would have said)

 Minor edit by pyraminx

 Minor edit by wisewigglyjaguar

 Minor edit by jeffyang2025

 Minor edit by Loading...

 



\textbf{Solution 2 (simpler)}

Draw segment $AE.$ Segment $AE$ is the diagonal of rectangle $ABEF,$ and its diagonals have length $\sqrt{7^2+1^2}=\sqrt{50}=5\sqrt2.$ From right triangle $AED,$ we use pythagorean theorem to find $DE = 5.$

 Now, we see similar triangles $\triangle CDE$ and $\triangle CBA$. Let $CE = a,$ and $CD = b.$ We can find that $AC = 5-b,$ and $CB = 7-a.$ These triangles have a ratio of $\frac {AB}{DE} = \frac{1}{5}.$ So we get that $\frac {5-b}{a} = \frac{1}{5}.$ Cross multplying, we get $a =25-5b.$ And also $\frac{CB}{CD} = \frac{1}{5} = \frac{7-a}{b}.$ Cross multiplying gives $35-5a=b.$ Solving the system of equations, we find $a = 25/4,$ which means $CB = 7-25/4 = 3/4.$ $[ABC] = CB/2,$ which gives $\boxed{[ABC] = 3/8}.$




\textbf{Solution 3 (thorough)}

From the diagram, $\angle BCA$ and $\angle DCE$ are vertical angles and hence congruent. Additionally, $\angle B = \angle D = 90^\circ$, so we have by AA Similarity that $\triangle BCA \sim \triangle DCE$.  

 Let $BC = x$ so $EC = 7 - x$ and $AC = y$ so $CD = 5 - y$. Since the two triangles are similar, we have $\frac{BC}{AC} = \frac{CD}{EC}$. Plugging in the variables gives $\frac{x}{y} = \frac{5 - y}{7 - x}$.  

 Cross multiplying yields $(7 - x)(x) = (5 - y)(y) \implies 7x - x^2 = 5y - y^2 \implies 7x + (y^2 - x^2) = 5y$.  

 By applying the Pythagorean Theorem into $\triangle BCA$, we get $x^2 + 1 = y^2 \implies 1 = y^2 - x^2$.  

 Therefore, $y = \frac{7x + 1}{5}$, and plugging this back into $x^2 + 1 = y^2$:

 $x^2 + 1= (\frac{7x+1}{5})^2$
$25(x^2 + 1) = (7x+1)^2$
$25x^2 + 25 = 49x^2 + 14x + 1$
$0 = 49x^2 + 14x + 1 - 25x^2 - 25$
$0 = 24x^2 + 14x - 24$
$0 = 12x^2 + 7x - 12$
$x = \frac{-7 \pm \sqrt{7^2 - 4(12)(-12)}}{2\cdot 12}$
$= \frac{-7 \pm \sqrt{49 + 576}}{24}$
$= \frac{-7 \pm \sqrt{625}}{24}$
$= \frac{-7 \pm 25}{24}$

 Therefore,$x = \frac{-7+25}{24}=\frac{18}{24}=\frac{3}{4}$.

 The area of $\triangle BCA$ is therefore $\frac{x \cdot 1}{2}=\frac{3}{8}=\boxed{A}$. ~hxve ~brilliantcoyote

 



\textbf{Solution 4 (Trigonometry)}

Using the Pythagorean theorem, I can get $AE=5 \sqrt{2}$. Then, because $AD=5$, $ED=5$. Now, let $\angle FEA=a$ and $\angle AED=b$. $\sin a=\frac{7}{5\sqrt{2}}, \sin b=\frac{1}{\sqrt{2}}, \cos a=\frac{1}{5\sqrt{2}},$ and $\cos b=\frac{1}{\sqrt{2}}$. Then, applying the sine addition formula, I get:

 $\sin(a+b)=\frac{7}{5\sqrt{2}}\times\frac{1}{\sqrt{2}}+\frac{1}{\sqrt{2}}\times\frac{1}{5\sqrt{2}}$
$=\frac{7}{10}+ \frac{1}{10}$
$=\frac{4}{5}$

 Thus, $\sin(a+b)=\frac{4}{5}$, so $\sin(180-a-b)=\frac{4}{5}$. This indicates that if a perpendicular is dropped from point $F$ to the extension of line $DE$, and the foot of the latitude is point $G$, $\triangle EGF$ is a 3-4-5 triangle. Because $\triangle EGF\sim\triangle CDE$, $\triangle CDE$ is also a 3-4-5 triangle. Using ratios,:

 $CE=\frac{5}{4}\times 5$
$=\frac{25}{4}$

 Therefore, $CE=\frac{25}{4}$, so $BC=\frac{3}{4}$, so $\triangle ABC$ has area $\frac{3}{8}$, or $\boxed A$.

 P.S. Thank you to eqb5000 and i\_am\_not\_suk\_at\_math for pointing out and helping me fix typos.

 



\textbf{Solution 5 (risky, but it works!)}

Using a ruler (which is permitted during the exam), and assuming the diagram is to-scale, we can measure the physical lengths of $AB$ and $BC$, and determine the scale factor in order to calculate $BC$'s actual math length. In my specific case (potentially could vary), $AB$ was $1$ cm and $BC$ was between $0.7$ and $0.8$ cm. So, the scale with cm is $1:1$, and the length of $BC$ is around $0.75$, so the area is $\frac{1}{2} \cdot 0.75 \cdot 1 = \frac{3}{8}$. To assure ourselves that $\frac{3}{8}$ is the most accurate estimation, we know that $\frac{4}{9}$ is around $0.44$ (too big), $\frac{\sqrt{13}}{8}$ even bigger (so also too big), $\frac{7}{15}$ is just under $0.5$, and $\frac{\sqrt{15}}{8}$ is even bigger than $\frac{\sqrt{13}}{8}$, so most likely the answer is $\boxed{\textbf{(A)} \frac{3}{8}}$.

 Remark - You could also redraw the figure to scale. That way, you won't have to assume it's drawn to scale in the provided figure.You would have to draw the diagonal AE, measure it as $\sqrt{7^2+1^2}$ and use pythagorean theorem to find ED as 5. Once you have that information you could sketch the whole figure to scale.

 



\textbf{Solution 6 (another cheese)}

Inspired by \href{https://artofproblemsolving.com/wiki/index.php?title=2022\_AMC\_10B\_Problems/Problem\_16}{2022 AMC 10B Problem 16}, another right-triangle and rectangle mashup, we can pretend that the problem uses $3-4-5$ similarity, because AMC problems do that a lot. In this case, we can verify it works by setting $BC = \frac{3}{4}, AC = \frac{5}{4}$ by $3-4-5$ similarity, so $CD = \frac{15}{4}$ and $CE = \frac{25}{4}$. Then, we see that $BC + CE = \frac{25+3}{4} = 7 = BE = AF$, which is our only other condition, so this setup works and we can bubble in $\frac{1}{2} \cdot \frac{3}{4} \cdot 1 = \boxed{\textbf{(A)} \frac{3}{8}}$. of course, if it weren't $3-4-5$, we can use any other technique.

 



\textbf{Solution 7 (area)}

Connect $AE$, which equals $\sqrt{50}$. Then calculate $ED$, which equals $5$.  $\triangle ABC \sim \triangle EDC$ (AA).  $\dfrac{AB}{ED} = \dfrac{1}{5}$.  

 Therefore,  $\dfrac{\text{Area of } \triangle ABC}{\text{Area of } \triangle EDC} = \dfrac{1}{25}$.  

 Also, let the area of $\triangle ACE = x$.  The area of $\triangle ABC = \dfrac{7}{2} - x$.  The area of $\triangle EDC = \dfrac{25}{2} - x$.  

 Solve:  $x = \dfrac{25}{8}$ 

 So,  $\text{Area of } \triangle ABC = \dfrac{7}{2} - \dfrac{25}{8} = \boxed{\textbf{(A)} \dfrac{3}{8}}$.

 



\textbf{Solution 8 (analytic geometry)}


                    \begin{figure}[H]
                        \centering
                        \includegraphics[width=0.5\textwidth]{Solution/AMC_10_10A_2025_15_Solution_Figure_1.png}
                    \end{figure}
                    As shown in the figure, establish a coordinate system with point A as the origin. 

 Suppose $C = (a, 1)$, then the equation of $AC$ is $y = \dfrac{1}{a}x$. Since $DE \perp AC$, the slope of $DE$ is $-a$, the equation of $DE$ is $y - 1 = -a(x - 7) \to ax + y - 7a - 1 = 0$. 

 Notice that the length of $AD$ is the distance between point $A$ and line $DE$, due to the distance formula, we get:
 \[\frac{|a \cdot 0 + 0 - 7a - 1|}{\sqrt{a^2 + 1}} = 5\]
 \[(7a + 1)^2 = 25(a^2 + 1)\]
 \[(4a - 3)(3a + 4) = 0\]
 \[a = \frac{3}{4}\]

 The area of $\triangle BCA$ is therefore $\frac{a \cdot 1}{2}=\frac{3}{8}=\boxed{A}$.

 



\textbf{Solution 9}

From the Pythagorean Theorem, $AE=5\sqrt{2}$ and we also discover $\triangle ADE$ is a $45°-45°-90°$ right triangle. This means that $DE=5$. As $\triangle ABC \sim \triangle EDC$ by AA similarity and $AB=1$, we can set up a system of linear equations in terms of $a$ and $b$ where they represent $BC$ and $AC$ respectively. Since $CD=5CB$ and $CE=5AC$, we have
 \[5a+b=5\]
 \[a+5b=7\]and solving the equations results in $a=\frac{3}{4}$. The area of $\triangle ABC$ is $\boxed{\textbf{(A)} \frac{3}{8}}$. $b$ is unnecessary to find.

 Note that $\triangle ABC$ is a $3-4-5$ triangle.

 This is very similar to Solution 2, and at the time of writing, I had not noticed.

 



\textbf{Solution 10}

Start by assigning the variable $x$ to the length of side $\overline{AC}$. Since $\overline{AD}=5$, $\overline {CD}=5-x$. Additionally, the diagonal $\overline{AE}=5\sqrt{2}$ by the Pythagorean Theorem. Notice that with this diagonal, we've created a right triangle $\triangle ADE$.

 

                    \begin{figure}[H]
                        \centering
                        \includegraphics[width=0.5\textwidth]{Solution/AMC_10_10A_2025_15_Solution_Figure_2.png}
                    \end{figure}
                    

 Now, you could manually use the Pythagorean Theorem, but you can just spot the $45^\circ-45^\circ-90^\circ$ triangle and see $\overline{ED}=5$.  After this, we can use the Pythagorean Theorem to find the length of $\overline{CE}$. 
 \[5^2+(5-x)^2   = (\overline{CE})^2\] 
 \[25+25+x^2-10x = (\overline{CE})^2\] 
 \[x^2-10x+50 = (\overline{CE})^2\] 
 \[\sqrt{x^2-10x+50}=\overline{CE}\]From this, we can see that $\overline{CB} = 7-\overline{CE} \implies \overline{CB} = 7 - \sqrt{x^2-10x+50}$
                    \begin{figure}[H]
                        \centering
                        \includegraphics[width=0.5\textwidth]{Solution/AMC_10_10A_2025_15_Solution_Figure_3.png}
                    \end{figure}
                    Now, we know that $\left[CDE\right] + \left[ABE\right] -\left[ABC\right] = \left[ ADE\right]= 12.5$. Therefore,
 \[\frac{5\cdot(5-x)}{2} + \frac{7}{2} - \frac{7 - \sqrt{x^2-10x+50}}{2}=12.5\]
 \[5\cdot(5-x) + 7 - (7 - \sqrt{x^2-10x+50})=25\]
 \[25-5x +  \sqrt{x^2-10x+50}=25\]
 \[\sqrt{x^2-10x+50}=5x\]
 \[x^2-10x+50=25x^2\]
 \[0=24x^2+10x-50\]
 \[0=12x^2+5x-25\]
 \[0=(4x-5)(3x+5)\]Now, we know that $x$ must be positive so our only valid solution is
 \[x=\frac{5}{4}\]Finally, we can substitute this value of $x$ into our expression for length $\overline {CB}$.  The other leg would just be $1$ making for a total area of $\frac{1}{2} \cdot\overline{CB}$. Hence,
 \[\left[ ABC\right]= \frac{7 - \sqrt{x^2-10x+50}}{2}\]
 \[\left[ ABC\right]= \frac{7 - \sqrt{\frac{5}{4}^2-10\cdot \frac{5}{4}+50}}{2}\]
 \[\left[ ABC\right]= \frac{7 - \sqrt{\frac{25}{16}-\frac{50}{4}+\frac{200}{4}}}{2}\]
 \[\left[ ABC\right]= \frac{7 - \sqrt{\frac{25}{16}+\frac{150}{4}}}{2}\]
 \[\left[ ABC\right]= \frac{7 - \sqrt{\frac{25}{16}+\frac{600}{16}}}{2}\]
 \[\left[ ABC\right]= \frac{7 - \sqrt{\frac{625}{16}}}{2}\]
 \[\left[ ABC\right]= \frac{7 - \frac{25}{4}}{2}\]
 \[\left[ ABC\right]= \frac{\frac{28}{4} - \frac{25}{4}}{2}\]
 \[\left[ ABC\right]= \frac{\frac{3}{4}}{2}\]
 \[\left[ ABC\right]=\frac{3}{8} \implies \boxed{\textbf{(A)} \frac{3}{8}}\]~Shockfront99

 



\textbf{Solution 11}

Connect $AE$, we can find that $AE=\sqrt{AB^2+BE^2}=\sqrt{1^2+7^2}=5\sqrt{2}$.And $\angle ADE=90^\circ$, $AE=5\sqrt{2}$, $AD=5$, then we can get $ED=5$, $\triangle AED$ is an isosceles right triangle 

 Then, $\angle AED=45^\circ$Let $\angle CED=\angle \alpha$, $\angle \alpha=\angle AED-\angle AEC$

 Because of the tangent subtraction formula $\tan(\alpha-\beta)=\frac{\tan\alpha-\tan\beta}{1+\tan\alpha\tan\beta}$, we can get$\tan\alpha=\frac{\tan\angle AED-\tan\angle AEC}{1+\tan\angle AED\cdot\tan\angle AEC}$

 Since $\tan\angle AED=\tan 45^\circ=1$, and $\tan\angle AEC=\frac{AB}{EB}=\frac{1}{7}$, we can get$\tan\alpha=\frac{3}{4}$

 Because $\frac{CD}{ED}=\tan\alpha=\frac{3}{4}$, then $CD=\frac{15}{4}$, so $EC=\frac{25}{4}$.

 Thus, $BC=\frac{3}{4}$.

 So, the area of $\triangle ABC=\frac{1}{2}\cdot BC\cdot AB=\frac{1}{2}\cdot\frac{3}{4}=\frac{3}{8}=\boxed{A}$.

 



\textbf{Solution 12 (Similar to Solution 3)}

We notice that by vertical angles that $\angle DCE = \angle BCA$, and we also have $\angle CDE = \angle ABC = 90$. $\triangle BCA$ is similar to $\triangle DCE$ by Angle-Angle Similarity.Now let $BC = x$. Using the Pythagorean Theorem on $\triangle ABC$, we have $CA = \sqrt{x^2 + 1}$. Since $ABEF$ is a rectangle, we have $AF = 7 = BE$, which also means $CE = 7 - x$. Since $AD = 5$, $DC = 5 - \sqrt{x^2 + 1}$. From our similar triangles, we have $\frac{BC}{DC} = \frac{CA}{CE} \Longrightarrow \frac{x}{5 - \sqrt{x^2 + 1}} = \frac{\sqrt{x^2 + 1}}{7-x}$. Cross multiplying, simplifying, and squaring gives us $7x - x^2 = 5\sqrt{x^2 + 1} - x^2 - 1 \Longrightarrow 7x + 1 = 5\sqrt{x^2 + 1} \Longrightarrow 49x^2 + 14x + 1 = 25x^2 + 25 \Longrightarrow 12x^2 + 7x - 12 = 0$. We can use the quadratic formula, and take the positive solution. $x = \frac{-7 + \sqrt{7^2 - 4(12)(-12)}}{2(12)} = \frac{3}{4}$. Since ABC is a right triangle, the area is $\frac{1}{2} (AB)(BC) = \frac{1}{2} (1)(\frac{3}{4}) = \boxed{\textbf{(A)} \frac{3}{8}}$, which is our final answer.

 


	\item (\textit{AMC 10 10A 2025}) There are three jars. Each of three coins is placed in one of the three jars, chosen at random and independently of the placements of the other coins. What is the expected number of coins in a jar with the most coins?


 $\textbf{(A) } \frac{4}{3} \qquad\textbf{(B) } \frac{13}{9} \qquad\textbf{(C) } \frac{5}{3} \qquad\textbf{(D) } \frac{17}{9} \qquad\textbf{(E) } 2$

 \textbf{Answer: }D



\textbf{Solution 1 (Case Work)}

We have three coins and three jars. Each coin is placed independently and randomly into one of the jars. Let $M$ be the maximum number of coins in any jar. We want to compute the expected value of $M$.

 Step 1: Count total outcomes

 Each coin has $3$ choices, so the total number of equally likely placements is $3^3 = 27$.

 Step 2: Casework on the maximum number of coins

 Case 1:  $M = 1$. This occurs when each jar has exactly one coin. There are $3! = 6$ assignments of coins to jars. Hence, $\Pr(M=1) = \frac{6}{27} = \frac{2}{9}$.

 Case 2:  $M = 3$. This occurs when all three coins fall into the same jar. There are $3$ jars to choose from, so $\Pr(M=3) = \frac{3}{27} = \frac{1}{9}$.

 Case 3:  $M = 2$. This occurs when one jar has $2$ coins, another jar has $1$ coin, and the last jar has $0$ coins.  We can choose which jar gets $2$ coins in $3$ ways, which jar gets $1$ coin in $2$ ways, and which $2$ coins out of the $3$ go into the jar with two coins, so we multiply by $\dbinom{3}{2}$, which is just $3$ (note we don't have to do this for the earlier cases because for case $2$, all $3$ coins go into one jar, and for case $1$, the factorial already accounts for that). Therefore, there are $3^2 \cdot 2 = 18$ outcomes. Thus, $\Pr(M=2) = \frac{18}{27} = \frac{2}{3}$.

 Step 3: Compute the expected value The expected value of $M$ is $\mathbb{E}[M] = 1 \cdot \frac{2}{9} + 2 \cdot \frac{2}{3} + 3 \cdot \frac{1}{9}$.  Converting everything to ninths, we have $\mathbb{E}[M] = \frac{2}{9} + \frac{12}{9} + \frac{3}{9} = \frac{17}{9}$.

 Hence, the expected number of coins in the jar with the most coins is$\boxed{\text{(D) }\frac{17}{9}}$.

 $Pr$=Probability of

 $\mathbb{E}$=Expected value

 


 \textbf{Another way of finding cases of M = 2} As described in the solution, there are $3^3=27$ ways of distributing the coins into the $3$ jars. Because there are $6$ ways for M=1 and $3$ ways for M=3, there are $27 - 6 - 3 = 18$ ways for $M=2$.

 



\textbf{Solution 2 (Fast)}

Assuming all jars and coins are distinct, there are $27$ total outcomes. $3!=6$ of them distribute exactly, and thus a max of, $1$ per cup, $3$ ways to choose a jar to put all $3$ coins in for a max of $3$. This leaves $27-6-3=18$ for a max of $2$, so the expected value is $\frac{6 \cdot 1 + 18 \cdot 2 + 3 \cdot 3}{27} = \boxed{\text{(D) }\frac{17}{9}}$




\textbf{Solution 3 (Circular Counting - Expectation Cycles)}

In circular counting, expected value is determined as
 \[\mathbb{E}_C n(A) = \frac{C n(A)}{\mathbb{T} - C n(A) - n(A)}\]Where \( \mathbb{E}_C n(A) \) is the expected value, \( C n(A) \) is the circular gap theorem, \( \mathbb{T} \) is the total number of cases, and \( n(A) \) is the invalid cases.

 Arrange the three jars and three coins into a circle of 6 elements. We want to count all the ways in which no 2 or no 3 coins are next to one another. This can easily be done using the circular gap theorem, which states
 \[C n(A) = \frac{n}{k} \binom{n-k-1}{k-1}\]Where \( n \) is the number of elements in our circle and \( k \) is the number of non adjacencies, in which we wish to count \( n = 6 \) and \( k \in \{1,2,3\} \) (The reason we have only 1, 2, and 3 is because if we have \( n=6 \) and \( k \ge 4 \), we see that the combinatorial part of \( C n(A) \) becomes invalid).

 So we begin counting.

 For \( k=1 \), using the formula, we have \( \frac{6}{1} \binom{4}{0} \) which is 6.

 For \( k=2 \), using the formula, we have \( \frac{6}{2} \binom{3}{1} \), which is 9.

 For \( k=3 \), using the formula, we have \( \frac{6}{3} \binom{2}{2} \) which is 2.

 Thus we have \( C_6(k \in \{1,2,3\}) = 6+9+2 = 17 \).

 Now we calculate the denominator of \( \mathbb{E}_C n(A) \). We see the total number of cases is \( 3^3 = 27 \). So we can find \( \mathbb{T} - C n(A) \). We check for \( n(A) \), in which through common sense, we can see that if the triplet $(a,b,c)$ represents the number of coins given to each jar, the triplet $(1,1,1)$ doesn't satisfy the condition, as all jars have an equal number of coins. Therefore \( n(A)=1 \).

 Our answer is
 \[\mathbb{E}_C n(A) = \frac{C n(A)}{\mathbb{T} - C n(A) - n(A)}\]
 \[\mathbb{E}_C n(A) = \frac{17}{27-17-1}\]
 \[\mathbb{E}_C n(A) = \boxed{\text{(D) }\frac{17}{9}}\]




\textbf{Solution 4 (Weighted Probabilities)}

WLOG label the jars as $1,2,3$ and the coins as $a,b,c$. The probability for a given coin to land in a given jar is 
 \[\frac{1}{3}\]The number of ways for all three coins to land in the same jar is $3$, either all in jar $1$, all in jar $2$, or all in jar $3$. Next, the number of ways to have two coins in one jar and one in another in another is $3^2 \cdot 2 = 18$. Finally, since there are $3^3=27$ ways, the last weight is 
 \[27-18-3=6\] Thus 
 \[\frac{1}{27} \cdot 3 \cdot 3 + \frac{1}{27} \cdot 2 \cdot 18 + \frac{1}{27} \cdot 1 \cdot 6\]Thus we get 
 \[\boxed{\text{(D) }\frac{17}{9}}\]




\textbf{Solution 5}

Denote the expected number of coins in a jar with the most coins as $E$
\begin{verbatim}
Process   $\llcorner$-- The 1st coin must add 1 to $E$. Situation: (1, 0, 0)       ├-- The 2nd coin has $1/3$ chance to add 1 to $E$. Situation: (2, 0, 0)       $\mid$   $\llcorner$-- The 3rd coin has $1/3$ chance to add 1 to $E$. Situation: (3, 0, 0)       $\llcorner$-- The 2nd coin has $2/3$ chance to leave $E$ unchanged. Situation: (1, 1, 0)           $\llcorner$-- The 3rd coin has $2/3$ change to add 1 to $E$. Situation: (2, 1, 0)
\end{verbatim}

Hence, $E = 1 + \dfrac{1}{3} \times (1 + \dfrac{1}{3} \times 1) + \dfrac{2}{3} \times (0 + \dfrac{2}{3} \times 1) = \boxed{\text{(D) }\frac{17}{9}}$




\textbf{Solution 6}

Create a tree diagram and add up the probabilities.

 $P = \frac{3}{9} + \frac{2}{9} + \frac{2}{9} + \frac{2}{9} + \frac{2}{9} + \frac{2}{9} + \frac{2}{9} + \frac{1}{9} + \frac{1}{9} = \boxed{\text{(D)}\frac{17}{9}}$

 Note:(It does not matter which jar you pick at the head of the diagram so long a you stay consistent)~Rames\_Jong

 



\textbf{Solution 7(Similar to Solution 4)}

We have $3$ cases namely

 All $3$ coins are in exactly one jar.$2$ coins are in one jar and the last one is in one of the other two jars.Each jar has exactly $1$ coin.

 For the first case, we note that there are exactly $3$ ways all coins are in one of the three jars.For the second case, we note that we have $3$ ways to select the jar that are $2$ coins and $2$ ways to select the jar that has $1$ coin which results in $3 \cdot 2 = 6$ ways. However, there are $\binom{3}{2}$ ways to select $2$ out of the three coins to put in the jar that has $2$ coins. This results in $6 \cdot \binom{3}{2} = 18$ ways.Finally, the third case results in $3! = 6$ ways since there are $3$ ways to select one of the coins to put in jar $1$ and $2$ ways to select the remaining coins to put in jar $2$ and lastly only the last coin can go in jar $3$ which results in $1$ way. Obviously, this gives $3 \cdot 2 \cdot 1 = 3! = 6$ ways.Finally, we execute the expected value. The probability for each jar is obviously $(\frac{1}{3})^3 = \frac{1}{27}$. Remember to include the value of each expected value as well. Case 1 has a value of $3$, case 2 has a value of $2$, and case 1 has a value of $1$ since each of these values is exactly the maximum number of coins in each case. We finally have

 $\frac{1}{27} \cdot 3 \cdot 3 + \frac{1}{27} \cdot 18 \cdot 2 + \frac{1}{27} \cdot 6 \cdot 1 = \frac{1}{3} + \frac{4}{3} + \frac{2}{9} = \frac{5}{3} + \frac{2}{9} = \frac{15}{9} + \frac{2}{9} = \frac{17}{9}$ so the answer is $D$.

 



\textbf{Solution 8 (Cheese)}

(Use if you don't know how expected value works).

 Since we know that there has to be at least $1$ coin in the jar with the most coins, we can instantly rule out answer choice $C$, as $5/8 < 1$. We also know $E$ can't be the answer, as there are more cases to get $1$ as the most rather than $3$. Since we also know there a majority of cases of $2$, the answer will lean towards $2$, giving us answer choice $D$.~mugim

 


	\item (\textit{AMC 10 10A 2025}) Let $N$ be the unique positive integer such that dividing $273436$ by $N$ leaves a remainder of $16$ and dividing $272760$ by $N$ leaves a remainder of $15$. What is the tens digit of $N$?


 $\textbf{(A) } 0 \qquad\textbf{(B) } 1 \qquad\textbf{(C) } 2 \qquad\textbf{(D) } 3 \qquad\textbf{(E) } 4$

 \textbf{Answer: }E



\textbf{Solution 1}

The problem statement implies that $N$ divides both $273436-16=273420$ and $272760-15=272745$. We want to find $N > 16$ that satisfies both of these conditions. Hence, we can just find the greatest common divisor of the two numbers. $\gcd(273420,272745)=\gcd(675,272745)=\gcd(675,45)=\gcd(0,45)=45$ by the Euclidean Algorithm, so the answer is $\boxed{\text{(E) }4}.$

 Note: If an integer $a$ is congruent to an integer $b$ modulo a positive integer $c$, denoted by $a \equiv b \pmod{c}$, this means that $c$ divides the difference, $|a-b|$ (Euclidean Algorithm)

 minor edits by AD\_12 and mw2014

 



\textbf{Solution 2}

We are given that $273436 \equiv 16 \pmod N$ and $272760 \equiv 15 \pmod N$, from which we get $273420 \equiv 0 \pmod N$ and $272745 \equiv 0 \pmod N$.

 Substracting the two gives $273420 - 272745 = 675$, which is also divisible by $N$.

 Since $N$ must be greater than $16$, the only possible divisors of $675$ are $25, 27, 45, 75, 135, 225,$ and $675$. Checking which ones also divide $273436$ and $273420$ eliminates $25$ and $27$, but $273420 \equiv 0 \pmod {45}$.

 Since both $273420$ and $675$ are divisible by $45$, so is $273420-675=272745$. Thus, $N = 45$ works, and its tens digit is $\boxed{\text{(E) }4}$.

 



\textbf{Solution 3}

We get that:$273436 \equiv {16}\pmod{N}$ and$272760 \equiv {15}\pmod{N}$.

 So we also have that:$273420 \equiv {0}\pmod{N}$ and$272745 \equiv {0}\pmod{N}.$

 Notice that these are a multiple of $5$. Now, we subtract these numbers to get $675$.We see that $675 = 25 \cdot 27.$ There is a factor of $5$ and $9$.$273420$ and $272745$ are multiples of $9$ as well, so our answer is just 45. Now the ten's digit is $\boxed{\text{(E) }4}$.

 



\textbf{Solution 4}

We are given that $273436 \equiv 16 \pmod{N}$ and $272760 \equiv 15 \pmod{N}.$  Subtracting, we get $273436 - 272760 = 676 \equiv 1 \pmod{N}.$ 

 This means $N$ divides $675.$  Also, since $273436 \equiv 16 \pmod{N},$ we have $273420 \equiv 0 \pmod{N}.$ 

 Thus, $N$ divides both $273420$ and $675.$ Using the Euclidean Algorithm:  $273420 - 405 \times 675 = 45,$ so $\gcd(273420, 675) = 45.$ 

 Hence $N = 45,$ and the tens digit of $N$ is $\boxed{\text{(E) }4}.$



	\item (\textit{AMC 10 10A 2025}) The \textit{harmonic mean} of a collection of numbers is the reciprocal of the arithmetic mean of the reciprocals of the numbers in the collection. For example, the harmonic mean of 4, 4, and 5 is 


 
 \[\frac{1}{\frac{1}{3}(\frac{1}{4}+\frac{1}{4}+\frac{1}{5})}=\frac{30}{7}\]


 What is the harmonic mean of all the real roots of the 4050th degree polynomial


 
 \[\prod_{k=1}^{2025} (kx^2-4x-3) = (x^2-4x-3)(2x^2-4x-3)(3x^2-4x-3)\dots (2025x^2-4x-3) ?\]
$\textbf{(A) } -\frac{5}{3} \qquad\textbf{(B) } -\frac{3}{2} \qquad\textbf{(C) } -\frac{6}{5} \qquad\textbf{(D) } -\frac{5}{6} \qquad\textbf{(E) } -\frac{2}{3}$

 \textbf{Answer: }B



\textbf{Solution 1}

We will need to determine the sum of the reciprocals of the roots. To find the sum of the reciprocals of the roots $p,q$ of the quadratic $ax^2+bx+c$, we use Vieta's formulas. Recall that $p+q = -b/a$ and $pq = c/a$. Therefore, 
 \[\frac{1}{p} + \frac{1}{q} = \frac{p+q}{pq} = \frac{\frac{-b}{a}}{\frac{c}{a}} = \frac{-b}{a} \cdot \frac{a}{c} = \frac{-b}{c},\]which doesn't depend on $a$.

 The sum of the reciprocals of the roots of the quadratic $x^2-4x-3$ is $\frac{-(-4)}{-3} = -4/3.$ The same is true for every quadratic in the form $ax^2-4x-3$. The sum of all the reciprocals of the roots of $ax^2+bx+c$ is $2025 \cdot \left(-\frac{4}{3}\right).$

 Because we have $2025$ quadratics, there are $2 \cdot 2025 = 4050$ total roots. Our answer is $\frac{1}{\frac{1}{4050}\cdot \frac{-4\cdot 2025}{3}} =\boxed{(B)\ -\dfrac{3}{2}}$.

 



\textbf{Solution 2 (similar to solution 1 but with quadratic formula)}

We first find the general roots for quadratics in the form $kx^2-4x-3$. Using the quadratic formula we have \begin{align*}x&=\frac{4 \pm \sqrt{16+12k}}{2k} \\&= \frac{4\pm 2\sqrt{4+3k}}{2k} \\&= \frac{2+\sqrt{4+3k}}{k}, \frac{2-\sqrt{4+3k}}{k} \\\end{align*}
Since we are asked to add the reciprocals of all $4050$ roots in the harmonic mean, we will first add the general roots in terms of $k$. We have,\begin{align*}\frac{k}{2+\sqrt{4+3k}} + \frac{k}{2-\sqrt{4+3k}} &= \frac{k(2-\sqrt{4+3k}) + k(2+\sqrt{4+3k})}{2^2-(\sqrt{4+3k})^2} \\&= -\frac{4k}{3k}=-\frac{4}{3}. \\\end{align*}
Thus, each pair of roots add up to $-\frac{4}{3}$, and since there are $2025$ pairs of roots, the harmonic mean of the desired expression is \begin{align*}\frac{1}{\frac{1}{4050} \left (2025 \left (-\frac43 \right ) \right)} &= \frac{1}{\frac12 \left ( -\frac43 \right)} \\& = \boxed{-\frac32}, \boxed{B}. \\\end{align*}


 



\textbf{Solution 3 (Vieta's only)}

We are asked to find $\frac{4050}{\sum_{i=1}^{4050}\frac{1}{i_1}}$. By Vieta's, note that $r_1 \cdot r_2 \dots r_{4050} = -\frac{3^{2025}}{2025!} = k$ ($k$ is a constant). Then, note that we are asked to find $\frac{4050}{\sum_{i=1}^n \frac{\frac{k}{r_i}}{k}}$, and by Vieta's we get that $\sum_{i=1}^n \frac{k}{r_i} = \frac{2025 \cdot 4 \cdot 3^{2024}}{2025!}$, so substituting this in, we get 
 \[-\frac{4050}{\frac{\frac{2025 \cdot 4 \cdot 3^{2024}}{2025!}}{\frac{3^{2025}}{2025!}}} = -\frac{4050 \cdot 3^{2025}}{4050 \cdot 3^{2024} \cdot 2} = \boxed{-\frac{3}{2}}\].

 


 



\textbf{Solution 4 (Vieta's Formulas)}

We have: 
 \[\prod_{k = 1}^{2025} (kx^{2} - 4x - 3) = (x^{2} - 4x - 3)(2x^{2} - 4x - 3)(3x^{2} - 4x - 3) \cdots (2025x^{2} - 4x - 3)\]

 So we can analyze each part of this "snake":

 According to Vieta's Formulas:

 For the first one: $x^{2} - 4x - 3$:We have:
 \[\begin{cases}  x_1+x_2=\frac{4}{1}\\ x_1x_2=-\frac{3}{1} \end{cases}\]So we have: 
 \[\dfrac{1}{x_1}+\dfrac{1}{x_2}=\dfrac{x_1+x_2}{x_1x_2}=-\dfrac{4}{3}\]

 For the second one: $2x^{2} - 4x - 3$:We have:
 \[\begin{cases}  x_3+x_4=\frac{4}{2}\\ x_3x_4=-\frac{3}{2} \end{cases}\]So we have: 
 \[\dfrac{1}{x_3}+\dfrac{1}{x_4}=\dfrac{x_3+x_4}{x_3x_4}=-\dfrac{4}{3}\]

 For the third one: $3x^{2} - 4x - 3$:We have:
 \[\begin{cases}  x_5+x_6=\frac{4}{3}\\ x_5x_6=-\frac{3}{3} \end{cases}\]So we have: 
 \[\dfrac{1}{x_5}+\dfrac{1}{x_6}=\dfrac{x_5+x_6}{x_5x_6}=-\dfrac{4}{3}\]

 ......

 Then we should check the last one. For the 2025th. one: $2025x^{2} - 4x - 3$:We have:
 \[\begin{cases}  x_{4049}+x_{4050}=\frac{4}{2025}\\ x_{4049}x_{4050}=-\frac{3}{2025} \end{cases}\]So we have: 
 \[\dfrac{1}{x_{4049}}+\dfrac{1}{x_{4050}}=\dfrac{x_{4049}+x_{4050}}{x_{4049}x_{4050}}=-\dfrac{4}{3}\]
 \begin{align*} \therefore H.M.&=\dfrac{1}{\dfrac{1}{4050}\sum_{i=1}^{4050}\dfrac{1}{x_i}}\\ &=\dfrac{4050}{2025\times\left(-\dfrac{4}{3}\right)}\\ &=-\dfrac{3}{2} \end{align*}
Therefore, the final answer is $\color{red}\boxed{\color{black}-\dfrac{3}{2}}\color{black}$.

 



\textbf{Solution 5 (Vietas Formulas)}

For a general quadratic polynomial $ax^2 + bx + c = 0$ with roots $p$ and $q$, the sum of the reciprocals is:\begin{align*}\frac{1}{p} + \frac{1}{q} &= \frac{p+q}{pq} \\\end{align*}
Using Vieta's formulas ($p+q = -\frac{b}{a}$ and $pq = \frac{c}{a}$):\begin{align*}\frac{1}{p} + \frac{1}{q} &= \frac{-\frac{b}{a}}{\frac{c}{a}} = -\frac{b}{c}\end{align*}
This result demonstrates that the sum of the reciprocals depends only on the coefficients $b$ and $c$, and is independent of the leading coefficient $a$. Since the overall polynomial is a product of quadratics, and all these quadratics share the same $b$ and $c$ coefficients, the sum of the reciprocals for the roots of every individual quadratic is the same. We use the coefficients from one of the quadratics, which is $ax^2 - 4x - 3 = 0$. Thus, $b=-4$ and $c=-3$.

 
 \[\sum (\text{reciprocals}) = -\frac{b}{c} = -\frac{(-4)}{(-3)} = -\frac{4}{3}\]

 We substitute the calculated sum of the reciprocals into the formula for the harmonic mean of 2 numbers:

 
 \[\text{HM} = \frac{2}{-\frac{4}{3}} = 2 \times \left(-\frac{3}{4}\right) = -\frac{6}{4} = -\frac{3}{2}\]

 The harmonic mean of the roots of any single quadratic is $-\frac{3}{2}$. Since the sum of the reciprocals is the same for all quadratics, the harmonic mean of all the roots of the overall polynomial is also $\boxed{\textbf{(B) }\dfrac{-3}{2}}$.

 \textbf{Note} It is important to note that the question asks for the sum of all \textit{real} roots. We must therefore be careful to make sure that the $4050$ roots of the polynomial are all real, and also that they are distinct, as otherwise e.g. double roots would be counted twice, when the problem requires us to only count each root once (i.e. regardless of multiplicity). Accordingly, we compute that each quadratic factor $\left(kx^2-4x-3\right)$ (for $1 \leq k \leq 2025$) has discriminant $(-4)^2-4 \cdot k \cdot (-3) = 16+12k$, which is clearly positive as $k \geq 1$. This means the $2$ roots of each quadratic factor are indeed real and distinct (i.e. none of these quadratics has a double root), so it only remains to show that these roots are also distinct from those of all the \textit{other} quadratic factors. Hence, if we assume that $a$ is a common root of $2$ different quadratic factors simultaneously, say $\left(px^2-4x-3\right)$ and $\left(qx^2-4x-3\right)$ with $1 \leq p,q \leq 2025$ and $p \neq q$, then we must have $pa^2-4a-3=qa^2-4a-3=0$, which implies $pa^2 = qa^2$. This equation factorises as $(p-q)a^2 = 0$, and because $p \neq q$, i.e. $p-q \neq 0$, it further reduces to simply $a^2 = 0$, and hence $a = 0$. But then it follows that $pa^2-4a-3 = p \cdot 0^2 - 4 \cdot 0 - 3 = -3 \neq 0$, yielding a contradiction, so $a$ cannot in fact have been a common root of $2$ different quadratic factors, and thus we finally deduce that all $4050$ roots are distinct, as required.~ Shadowleafy

 



\textbf{Solution 6}

Let the polynomial be $f(x),$ and denote the $4050$ roots to be $x_1,x_2,...,x_{4050}.$ Hence, 
 \[HM = \dfrac{4050}{\frac{1}{x_1}+\frac{1}{x_2}+...\frac{1}{x_{4050}}}.\]We can multiply the numerator and denominator of this fraction by $x_1x_2...x_{4050}$ to create symmetric sums, which yields 
 \[HM = \dfrac{4050(x_1x_2...x_{4050})}{x_2x_3...x_{4050}+x_1x_3...x_{4050}+...+x_1x_2...x_{4049}}.\]

 By Vieta's Formulas, since $f(x)$ is of even degree, the product of its roots, $x_1x_2...x_{4050},$ is just the constant term of $f(x),$ call it $c_0.$ Likewise, the denominator of our harmonic mean, $x_2x_3...x_{4050}+x_1x_3...x_{4050}+...+x_1x_2...x_{4049},$ is the negated coefficient of $x$ in the standard form of $f(x).$ Let the coefficient of $x$ in the standard form of $f(x)$ be $c_1.$ Note that we do not have to worry about dividing by the coefficient of $x^{4050}$ when using Vieta's Formulas, as they will eventually cancel out in our harmonic mean calculations.

 So, 
 \[HM = \dfrac{4050c_0}{-c_1}.\]

 The constant term in $f(x)$ is just $c_0=(-3)^{2025}.$ For the coefficient of the $x$ term in $f(x),$ there are $\dbinom{2025}{2024}=2025$ ways to choose $2024$ of the trinomials to include a $-3,$ and the one trinomial not chosen will include a $-4x.$ Hence, $c_1=2025\cdot (-3)^{2024}\cdot (-4).$

 Finally, 
 \[HM=\dfrac{4050\cdot(-3)^{2025}}{-2025\cdot (-3)^{2024}\cdot (-4)}=\dfrac{2\cdot(-3)}{-(-4)}=\boxed{\text{(B) }-\dfrac{3}{2}}.\]




\textbf{Solution 7 (Cheese )}

We can try replacing the $2025$ in the problem with smaller numbers like $1$ and $2$. For $1$, we have the quadratic $x^2-4x-3$. Using the quadratic formula, we have the roots are $2+\sqrt{7}$ and $2-\sqrt{7}$ The harmonic mean of these two is 
 \[\frac{2}{\frac{1}{2+\sqrt{7}} + \frac{1}{2-\sqrt{7}}}=\frac{2}{\frac{\sqrt{7}-2}{3}+\frac{2+\sqrt{7}}{-3}}=\frac{2}{\frac{-4}{3}}=-\frac{3}{2}\].

 We notice that this is one of the answer choices. Also, given that the random choice of $2025$, and that the rest of the answer choices are also simple fractions, we can safely guess that the answer is $\boxed{\textbf{(B)}~-\frac{3}{2}}$




\textbf{Solution 8 (Symmetry of Unweighted Means)}

Note that all the roots are real since for each quadratic, 
 \[\Delta=b^2-4ac=(-4)^2-4k(-3)=16+12k>0\] for $k\ge 1$. We will use the fundamental theorem of symmetric polynomials. By Vieta's formulas, we have for each quadratic $ax^2+bx+c=kx^2-4x-3$:
 \[\frac{1}{r}+\frac{1}{s}=\frac{r+s}{rs}=\frac{-b/a}{c/a}=-\frac{b}{c}=-\frac{4}{3}\]Thus, $\mathrm{HM}(r,s)=\frac{2}{\frac{1}{r}+\frac{1}{s}}=\frac{2}{-\frac{4}{3}}=-\frac{3}{2}$. Since the unweighted harmonic mean is symmetric in its arguments, we can group terms to get
 \[\mathrm{HM}(r_{1},r_{2},r_{3},\dots,r_{4050})=\mathrm{HM}\bigg(\mathrm{HM}(r_{1},r_{2}),\mathrm{HM}(r_{3},r_{4}),\dots,\mathrm{HM(r_{4049},r_{4050})}\bigg)=\boxed{\textbf{(B)}~-\frac{3}{2}}\]

 Note: For an unweighted symmetric mean $M$, we have for $d|n$, $\mathrm{M}\{ x_{k} \}_{k=1}^{n}=\mathrm{M}\{ \mathrm{M}\{ x_{ds+k} \}_{k=1}^{d} \}_{s=0}^{n/d-1}$ and $\mathrm{M}\{ \bar{x} \}_{k}=\bar{x}$ for constant mean $\bar{x}$. You can also check this individually for HM: 
 \[\mathrm{HM}(x_{1},\dots,x_{n})=n(x_{1}^{-1}+\cdots+x_{n}^{-1})^{-1}=\frac{n}{d}\left[ \underbrace{ d\left(x_{1}^{-1}+\cdots+x_{d}^{-1}\right)+\cdots+d\left(x_{(n-1)d+1}^{-1}+\cdots+x_{n}^{-1}\right) }_{ \frac{n}{d}\text{ terms} } \right]\]or by checking arithmetic mean and using $\mathrm{HM}\{ x_{k} \}_{k}=\frac{1}{\mathrm{AM}\left\{ \frac{1}{x_{k}} \right\}_{k}}$.

 



\textbf{Solution 9 (Easy-ish to understand)}

Note: I will try my best to make this solution look less intimidating to read as the other ones. This is my third time writing a solution, so I'm not very used to the $\LaTeX$. 

 The harmonic mean that the question is asking is 
 \[\frac{1}{\frac{r_1+r_2+...r_{4049}+r_{4050}}{4050}}\] where $r_i$ is between $1$ and $4050$ and is the reciprocal of all $4050$ roots of the polynomial. All we need to find is the sum of the reciprocals of the roots and we can substitute that into this formula. 

 First of all, notice that there are $2025$ quadratic polynomials. They each hold $2$ out of the $4050$ roots, and the sum of the reciprocals of their roots are $\frac{1}{a} + \frac{1}{b} = \frac{a+b}{ab}.$ 

 Using Vieta's we get $\frac{a+b}{ab}=-\frac{4}{3}$ on the every quadratic because only the coefficient for $x^2$ changes. Also, $a+b$ and $ab$ have the denominator of $a$, so when they are divided, the quotient won't change. Therefore, the sums of the reciprocals of the quadratics sum to $2025 \cdot \frac{-4}{3}$.

 Now we can substitute: 
 \[\frac{1}{\frac{2025 \cdot \frac{-4}{3}}{4050}}.\] This gets us \begin{align}\frac{1}{\frac{2025 \cdot \frac{-4}{3}}{4050}} &= \frac{1}{\frac{2025 \cdot (-4)}{3 \cdot 4050}} \\&= \frac{1}{\frac{-8100}{12150}} \\&= \frac{1}{\frac{-2}{3}} \\&= 1 \cdot \frac{3}{-2} \\&= -\frac{3}{2}\end{align}

 Or...you could just solve the last part yourself easier than my solution!

 



\textbf{Solution 10 (reciprocal roots)}

We know that if for polynomial $ax^2 + bx + c$ is $r_1$ and $r_2$, then the polynomial with roots $\frac{1}{r_1}$ and $\frac{1}{r_2}$ (reciprocal roots) will be $cx^2 + bx + a$ where the coefficients are flipped (works for any degree polynomial, not just quadratic). Thus, using vieta's formula, the sum of the reciprocal roots for polynomial $kx^2 - 4x - 3$ will be the sum of the roots for polynomial $-3x^2 - 4x + k$ which is $- \frac{-4}{-3}=- \frac{4}{3}$.

 There are $2025$ polynomials like this, so the sum of roots would be $2025\cdot (- \frac{4}{3})$. Finding the harmonic mean, we get:

 
 \[\frac{1}{\frac{1}{4050} \cdot \frac{-4\cdot 2025}{3}} = \frac{1}{\frac{-2}{3}}\]

 Or $\boxed{\textbf{(B) }-\frac{3}{2}}$



	\item (\textit{AMC 10 10A 2025}) An array of numbers is constructed beginning with the numbers 
$-1\qquad3\qquad1$in the top row. Each adjacent pair of numbers is summed to produce a number in the next row. Each row begins and ends with 
$-1$ and $1$, respectively.


 
 \[\begin{tabular}{ccccccccc} &&-1&&3&&1&&\\ &-1&&2&&4&&1&\\ -1&&1&&6&&5&&1\\ \end{tabular}\]


 If the process continues, one of the rows will sum to $12,288$. In that row, what is the third number from the left?


 $\textbf{(A) } -29 \qquad\textbf{(B) } -21 \qquad\textbf{(C) } -14 \qquad\textbf{(D) } -8 \qquad\textbf{(E) } -3$

 \textbf{Answer: }A



\textbf{Solution 1}

Consider the polynomial $f(x) = -x^2+3x+1.$ When we multiply this polynomial by $x+1,$ we are essentially doing the operation given in the problem (When we multiply $p(x)$ by $x+1,$ a term of degree $d$ in the yielded expression is the sum of $1\cdot(\text{degree d})$ and $x\cdot(\text{degree d-1})$ in $p(x)$ This effect is visible in Pascal's Triangle). So, if we let the coefficients of $f(x)$ be the zero row of the array, then the $n^{th}$ row is just the coefficients of $f(x)(x+1)^n.$The next thing to note is that the sum of the coefficients in any polynomial $p(x)$ is just $p(1).$ Therefore, the sum of the entries in the $n^{th}$ row of the array is $f(1)(1+1)^n=3\cdot2^n.$ Letting this equal $12288,$ we get $n=12.$ We are looking for the $3^{rd}$ term in the $12^{th}$ row. The $12^{th}$ row is given by the coefficients of $f(x)(x+1)^{12}=(-x^2+3x+1)(x+1)^{12}.$ Since the degree of the resulting expression is $14,$ the third term in the row is just the coefficient of $x^{12}$ in the expression, which is $-\dbinom{12}{10}+3\dbinom{12}{11}+1=\boxed{\textbf{(A) }-29}.$




\textbf{Solution 2}

If we take a look at the first few rows, we notice that the sum of the terms in each row $n$ is equal to the twice the sum of row $n-1$. We note the first row is $3$ so recognize $12,228$ must be equal to $3$ times a power of $2$. $12,288=3\cdot 4096=3 \cdot 2^{12}$. Therefore, we are looking for the $3$rd term from the left in the $13$th row. From here, we can take a look at the second number of each row, noting that it will always decrease by $1$ because it is being added to the $-1$ on its left. We can compute the $3$rd term of each row by adding the $2$nd and $3$rd term of the previous row, which means that to find the $3$rd term of the nth row, we simply have to add all of the 2nd terms of each row starting from the row with the first appearance of the $3$rd term, which is $6$. Using this, we get the expression $6+(1+0+(-1)+...+(-(n-5))$, where $n$ is $13$. Plugging it in, we get $6+(1+0+(-1)+...+(-8))$, which tells us that the $3$rd number from the left of the $13$th row is $\boxed{\text{(A) -29}}$.

 


 \textbf{Remark (Why each row is twice the previous row)} To prove that the sum of each row is $2$ times the previous row, use the equation ${(1+1)}^n$ with the normal Pascal's triangle and use the binomial theorem from there.

 Example: In the fifth row of Pascal's triangle, the numbers are $1, 4, 6, 4, 1$. The first number can be represented by ${4 \choose 0} = 1$, the second by ${4 \choose 1} = 4$, the third by  ${4 \choose 2} = 6$, the fourth by  ${4 \choose 3} = 4$, and the last by  ${4 \choose 4} = 1$. The sum of these five numbers is $16$ which is also $2^{4}$. Using the binomial theorem with ${(1+1)}^4$, we get $1 + 4 + 6 + 4 + 1 = 16$, the same results in Pascal's triangle. We could also solve ${(1+1)}^4$ by converting it to $2^4 = 16$, proving the identity that all rows in Pascal's triangle are exactly twice the previous row and can be represented by $2^{n}$.

 



\textbf{Solution 3 - Fast  and Effective!}

\(\large \textbf{<3 minute solve!}\)

 Add all the numbers up in the first row. You get $3.$ Now, add all the numbers up on the second row. You get $6.$ Notice that as the rows keep going, the sum of the numbers in the row keeps doubling. When you repeat this process, you realize that you reach $12288$ on the $13^{th}$ row.

 Knowing this, we can use this pattern to find the solution quickly. We know that the first number will always be $-1$, so we can ignore that. Knowing this, all you have to do now is to add up the second and third numbers $13$ times. This is already done for us three times, so we just have to do it ten more. Note that you do not have to do this process for all the numbers, only the second and third.

 


 Doing this $13$ times (since $\frac{12288}{3}=4096=2^{12}=2^{13-1}$), you get the following string of numbers (starting from the first one): $1, 4, 6, 7, 7, 6, 4, 1, -3, -8, -14, -21,$ and finally, $\boxed{\textbf{(A) }-29}.$




\textbf{Solution 4}

From adding the sum of all terms in the first few rows, we see that $S(r(1))=3, S(r(2))=6, S(r(3))=12, S(r(4))=24,...$ and so on. We can deduce that $S(r(n))=3 \cdot 2^{n-1}$ is the sum of all numbers in row $n$. Now, set $3 \cdot 2^{n-1} = 12288$ so we have $2^{n-1} = 4096 = 2^{12}$. It follows that $n=13$. Now let $a_n$ be the $2^{nd}$ number of each row, and $b_n$ be the $3^{rd}$ number of each row. Since, the first number of each row is -1, $a_n=a_{n-1}-1$. Additionally, $b_n = b_{n-1}+a_{n-1}$.

 \begin{array}{|c|c|}\hlinea_n & b_n \\\hlinea_1 = 3 & b_1 = 1 \\a_2 = 2 & b_2 = 4 \\a_3 = 1 & b_3 = 6 \\a_4 = 0 & b_4 = 7 \\a_5 = -1 & b_5 = 7 \\a_6 = -2 & b_6 = 6 \\a_7 = -3 & b_7 = 4 \\a_8 = -4 & b_8 = 1 \\a_9 = -5 & b_9 = -3 \\a_{10} = -6 & b_{10} = -8 \\a_{11} = -7 & b_{11} = -14 \\a_{12} = -8 & b_{12} = -21 \\a_{13} = -9 & b_{13} = -29 \\\hline\end{array}

 We are asked to compute $b_{13} = \boxed{\textbf{(A) }-29}.$




\textbf{Solution 5 (Wish I used this)}

So the sum is doubling every time, starting with 3.

 $12288 = 3 \cdot 2^{12}$

 So in the 13th row we go, so we eventually get that the 3rd number is $\boxed{\text{(A)}-29}$

 Go check out solution 1, looks interesting.

 



\textbf{Solution 6}

We firstly notice that the pattern for the sums of each row is $S_n$ $=$ $3 \cdot 2^n$ where n is the number of the row. Since the sum for the row we are looking for is $12288$, We can write the equation $3 \cdot 2^n$ $=$ $12288$ which solving gives $n = 12$.

 
Recognizing that the $n^{th}$ row term is basically a scaled version of the first row-term $(-1,3,1)$, we can write a polynomial      $f(x) = (-1x+3x+1x^2)(x+1)^n$, where $(x+1)^n$ is kind of a "shifting step" to proceed down the triangle to the $n^{th}$ term.

 
-you can test for the different values like $(x+1)^1$ which results to coefficients $(-1,2,4,1)$ which matches the second row down

 
Now we recognize we can write the polynomial of the $12^{th}$ term in the form $(a_0 + a_1 x + a_2 x^2 + \dots + a_m x^m)(1+x)^n$ which we can simplify using the binomial theorem achieving this: $\sum_{i=0}^{m} a_i \binom{n}{k-i}$ 

 
Plugging the terms in we get: $(-1)\binom{n}{2} + 3\binom{n}{1} + 1\binom{n}{0}$ where $n = 12$ thus giving$\boxed{(A) -29}$




\textbf{Solution 7}

The first step is recognizing that you multiply the sum by two each row. This is because each number adds to the number down to the left and the number down to the right (this includes the -1 and 1 on the ends). Thus the sum in row $n$ is $3*2^{n-1}$ which means the question corresponds to row $13$.

 Now we can look at the trends in the terms of each row. The first term in each row is $-$1, which makes the second term in each row $3,2,1,0,-1...$

 This leads us to finding the third term in each row has the pattern of $1, 1+3, 1+3+2, 1+3+2+1, 1+3+2+1+0, 1+3+2+1+0+-1...$

 Therefore we just need to find $1$ plus the sum of $3$ through $-8$ inclusive which by any number of methods computes to $1-30=\boxed{\text{(A) }-29}$.

 



\textbf{Solution 8}

We first observe the sums of the given rows:

 \begin{array}{|c|c|}\hlineRow & Sum \\\hline1 & 3 \\2 & 6 \\3 & 12 \\4 & 24 \\\hline\end{array}

 This sequence doubles each time, so the sum of the nth row is 
 \[3\cdot2^{n-1}.\] Setting this equal to $12,288$, we get\begin{align*}3\cdot 2^{n - 1} &= 12,288 \\2^{n - 1} &= 4096 \\4096 &= 2^{12} \\2^{n-1} &= 2^{12},\end{align*}
so $n=13$. Thus, the row with sum $12,288$ is the 13th row.Next, notice that every entry in the triangle below is formed by repeatedly summing from the three numbers in the first row: 1, 3, and -1. 

 
 \[\large{-1}\qquad\large{3}\qquad\large{1}\]
 \[\large{-1}\qquad\large{2}\qquad\large{4}\qquad\large{1}\]
 \[\large{-1}\qquad\large{1}\qquad\large{6}\qquad\large{5}\qquad\large{1}\]

 For example, the number $6$ in row $3$ can be expressed as\begin{align*}6 = 1 \cdot -1+2 \cdot 3+1 \cdot 1,\end{align*}
where the coefficients count the number of paths from each top entry. In general, each entry is a weighted sum of the top row, with weights equal to the number of left/right paths leading to that position.

 We are searching for the value of the 3rd entry of the 13th row. From the top row to this position, there are $12$ moves. To reach the 3rd position from:

 
            \begin{itemize}
        \item the 1st entry of row $1$, we need $2$ right moves and $10$ left moves.
            \end{itemize}
            Number of paths: $\dbinom{12}{2} = 66$.

 Contribution: $66 \cdot -1 = -66$. $\newline$

            \begin{itemize}
        \item the 2nd entry of row $1$, we need $1$ right move and $11$ left moves.
            \end{itemize}
            Number of paths: $\dbinom{12}{1}=12$.

 Contribution: $12 \cdot 3=36$. $\newline$

            \begin{itemize}
        \item the 3rd entry of row $1$, we need $0$ right moves.
            \end{itemize}
            Number of paths: $\dbinom{12}{0}=1$.

 Contribution: $1 \cdot 1=1$. $\newline$

 Adding the contributions gives\begin{align*}-66+36+1 = -29.\end{align*}
So, the 3rd number in the 13th row is $\boxed{\text{(A) }-29}$.

 


	\item (\textit{AMC 10 10A 2025}) A silo (right circular cylinder) with diameter 20 meters stands in a field. MacDonald is located 20 meters west and 15 meters south of the center of the silo. McGregor is located 20 meters east and $g > 0$ meters south of the center of the silo. The line of sight between MacDonald and McGregor is tangent to the silo. The value of $g$ can be written as $\frac{a\sqrt{b}-c}{d}$, where $a,b,c,$ and $d$ are positive integers, $b$ is not divisible by the square of any prime, and $d$ is relatively prime to the greatest common divisor of $a$ and $c$. What is the value of $a+b+c+d$?


 $\textbf{(A) } 119 \qquad \textbf{(B) } 120 \qquad \textbf{(C) } 121 \qquad \textbf{(D) } 122 \qquad \textbf{(E) } 123 \qquad$

 \textbf{Answer: }A



\textbf{Solution 1}


                    \begin{figure}[H]
                        \centering
                        \includegraphics[width=0.5\textwidth]{Solution/AMC_10_10A_2025_20_Solution_Figure_1.png}
                    \end{figure}
                    

 Let the silo center be $O$, let the point MacDonald is situated at be $A$, and let the point $20$ meters west of the silo center be $B$. $ABO$ is then a right triangle with side lengths $15, 20,$ and $25$. 

 
Let the point $20$ meters east of the silo center be $C$, and let the point McGregor is at be $D$ with $CD=g>0$. Also let $AD$ be tangent to circle $O$ at $E$.

 
Extend $BC$ and $AD$ to meet at point $F$. This creates $3$ similar triangles, $\triangle ABF\sim \triangle DCF \sim \triangle OEF$. Let the distance between point $C$ and $F$ be $x$. The similarity ratio between triangles $ABF$ and $DCF$ is then $\frac{longer\;leg}{shorter\;leg} = \frac{40+x}{15} = \frac{x}{g}$ 

 
This is currently unsolvable so we bring in triangle $OEF$. The hypotenuse of triangle $OEF$ is $OF=20+x$ and its shorter leg is the radius of the silo $=10$. We can then establish a second similarity relationship between triangles $OEF$ and $ABF$ with $\frac{shorter\; leg}{hypotenuse}=\frac{10}{20+x}=\frac{15}{AF}$

 
Now we find the hypotenuse of $ABF$ in terms of $x$ using the Pythagorean theorem. $AF^2=15^2+(40+x)^2$. Which simplifies to $AF^2=225+1600+80x+x^2=1825+80x+x^2$   So $AF=\sqrt{x^2+80x+1825}$

 
Plugging back in we get $\frac{10}{20+x}=\frac{15}{\sqrt{x^2+80x+1825}}$. Now we can begin to break this down by multiplying both sides by both denominators. $10(\sqrt{x^2+80x+1825})=15(20+x)$ Dividing both sides by $5$ then squaring yields, $4x^2+320x+7300=9x^2+360x+3600$  This furthermore simplifies to $5x^2+40x-3700=0$ At which point we can divide off a $5$ and then apply the quadratic formula on $x^2+8x-740=0$ which we take the positive root of.

 

 \[x = \frac{-8+\sqrt{64+2960}}{2} = \frac{-8+\sqrt{3024}}{2} =\frac{-8+\sqrt{144 \cdot 21}}{2}.\]

 
Simplifying yields that $x=6\sqrt{21}-4$

 
Then to solve for $g$ we simply plug $6\sqrt{21}-4$ back into the first similarity ratio to get $\frac{36+6\sqrt{21}}{15}=\frac{6\sqrt{21}-4}{g}$

 
Multiply both sides by $15g$ and dividing by $36+6\sqrt{21}$ will let us solve for $g=\frac{15(6\sqrt{21}-4)}{36+6\sqrt{21}}$ and after rationalizing the denominator we get $\frac{20\sqrt{21}-75}{3}$. $20+21+75+3=\boxed{\textbf{(A)}~119}$
$\sim$ Nioronean~ happyfish0922 (minor edit regarding what was previously point T changed to point E)~ Avs2010 (diagram)

 



\textbf{Solution 2 (Coordinates)}

Let MacDonald's position be at $(0,0)$. Then, the center of the silo will be at $(20,15)$ and McGregor will be at $(40, 15-g)$.

 Let the line of sight between MacDonald and McGregor be represented by the line $y=mx$, where $m=\frac{15-g}{40}$. As the radius of the silo is $10$, we can construct the equation of a circle and get the system of equations: 
 \[y=mx\] 
 \[(x-20)^2+(y-15)^2=100\]

 Substituting, we get: 
 \[(x-20)^2+(mx-15)^2=100\] 
 \[x^2-40x+400+m^2x^2-30mx+225=100\] 
 \[(m^2+1)x^2-(30m+40)x+525=0\]

 We want a value of $m$ such that this equation has $1$ solution, so we set the discriminant to be equal to $0$. 
 \[(30m+40)^2-4(m^2+1)(525)=0\] 
 \[900m^2+2400m+1600-2100m^2-2100=0\] 
 \[12m^2-24m+5=0\] 
 \[m=\frac{24 \pm \sqrt{576-240}}{24}\] 
 \[m=\frac{6 \pm \sqrt{21}}{6}\]

 The slope obviously has to be less than $1$, so we take the negative root $\frac{6-\sqrt{21}}{6}$.

 Substituting back for $m$ and solving for $g$, 
 \[\frac{15-g}{40}=\frac{6-\sqrt{21}}{6}\] 
 \[15-g=\frac{20}{3}(6-\sqrt{21})\] 
 \[g=\frac{20\sqrt{21}-75}{3}\]

 Finally, $a=20$, $b=21$, $c=75$, and $d=3$, summing up to a total of $20+21+75+3=\boxed{\textbf{(A)}~119}$




\textbf{Solution 3 (Trig)}

(I don't know how to make a diagram or latex very well oops)

 Let the silo center be $O$, let the point MacDonald is situated at be $A$, and let the point $15$ meters south of the silo center be $B$. $ABO$ is then a right triangle with side lengths $15, 20, 25$. 

 Also let the tangency point be $E$. Then, because a tangent line is defined as having a right angle to the radius $OE$, triangle $AEO$ is also a right triangle with side lengths $5\sqrt{21},10,25$ since we know $AO$ is 25 from pythag on $ABO$ and that $OE$ is 10 because it's a radius, so $AE$ follows from pythag on $AEO$.

 Let the position McGregor is in be $G$ and let the point that is both directly south of McGregor and directly east of MacDonald be $F$. (For reference, $F$ is $20$ meters east and $15$ meters south of O, which is also $40$ meters east of D and $15-g$ meters south of $G$.)

 Now, we can see that the triangle $AFG$ is also a right triangle. Since $G$ is on the line $AE$ by definition and $F$ is on line $AB$ because they are all on the same horizontal line $15$ meters south of O, $\angle GAF=\angle EAB$. This leads us to conclude that $\tan{\angle GAF}=\frac{15-g}{40}=\tan{\angle EAB}$.

 From here, we can use tangent subtraction formula to find $\tan(\angle EAB)$.

 We know that $\tan(\angle OAB)=\frac{15}{20}=\frac34$ and that $\tan(\angle OAE)=\frac{10}{5\sqrt{21}}\frac2{\sqrt{21}}$. Since $\angle EAB = \angle OAB - \angle OAE$, we have that

 $\tan(\angle EAB)=\tan(\angle OAB - \angle OAE)=\frac{\tan(\angle OAB)-\tan(\angle OAE)}{1+\tan(\angle OAB)\tan(\angle OAE)}=\frac{\frac34-\frac2{\sqrt{21}}}{1+\frac6{4\sqrt{21}}}$

 Multiplying the top and bottom by $4\sqrt{21}$, this equals

 $\frac{3\sqrt{21}-8}{4\sqrt{21}+6}=\frac12\cdot\frac{3\sqrt{21}-8}{2\sqrt{21}+3}$ (factoring out $\frac12$ is personal preference)

 Now, rationalizing the denominator by multiplying the top and bottom by $2\sqrt{21}-3$, we get this equals

 $\frac12\cdot\frac{(3\sqrt{21}-8)\cdot(2\sqrt{21}-3)}{(2\sqrt{21}+3)\cdot(2\sqrt{21}-3)}=\frac12\cdot\frac{150-25\sqrt{21}}{75}=\frac{6-\sqrt{21}}6$

 Recalling that this implies $\frac{15-g}{40}=\frac{6-\sqrt{21}}6$, we get that $g=\frac{20\sqrt{21}-75}{3}$. $20+21+75+3=\boxed{\textbf{(A)}~119}$




\textbf{Solution 4}


                    \begin{figure}[H]
                        \centering
                        \includegraphics[width=0.5\textwidth]{Solution/AMC_10_10A_2025_20_Solution_Figure_2.png}
                    \end{figure}
                    

 We use the point names given in Solution 1.

 By Pythagorean Theorem (twice), we have 
 \[AE=\sqrt{AO^2-OE^2}=\sqrt{AB^2+BO^2-10^2}=\sqrt{525}=5\sqrt{21}.\] Similarly, we find that 
 \[DE=\sqrt{DO^2-OE^2}=\sqrt{CD^2+CO^2-10^2}=\sqrt{300+g^2}.\] Note that the altitude from $D$ to $AB$ (call the foot $H$) has length $40$. Also, $AD=AE+DE=5\sqrt{21}+\sqrt{300+g^2}$ and $AH=15-g$. By Pythagoras on $\triangle ADH$, we have 
 \[\left(5\sqrt{21}+\sqrt{300+g^2}\right)^2=40^2+(15-g)^2.\] Expanding, 
 \[525+10\sqrt{21}\cdot\sqrt{300+g^2}+300+g^2=1600+225-30g+g^2\] and simplifying gives 
 \[10\sqrt{6300+21g^2}=1000-30g.\] Dividing by $10$ and squaring, 
 \[6300+21g^2=10000-600g+9g^2\] and putting it in standard form, 
 \[0=12g^2+600g-3700.\] Dividing by $4$ and applying the quadratic formula, 
 \[g=\frac{-150\pm\sqrt{22500-4(3)(-925)}}{6}=\frac{-75\pm20\sqrt{21}}3.\] From the answer form, we take the positive root, so 
 \[g=\frac{20\sqrt{21}-75}3.\] Finally, we have $20+21+75+3=\boxed{\text{(A) }119}.$




\textbf{Solution 5 (trig & similarity)}


                    \begin{figure}[H]
                        \centering
                        \includegraphics[width=0.5\textwidth]{Solution/AMC_10_10A_2025_20_Solution_Figure_2.png}
                    \end{figure}
                    

 First, taking the Pythagorean Theorem on $\Delta BOD$ and $\Delta AOE$ gives $AO = 25, AE = 5\sqrt{21}$. By cosine angular addition, $\cos{(\angle BOD + \angle AOE)} = \cos{\angle BOE} = \frac{8-3\sqrt21}{25}$, so $\cos{\angle EOC} = \frac{3\sqrt21-8}{25}$. 

 Dropping the altitude from $E$ to $BC$ and labelling the foot $P$, we have that $OP = 10 \cdot \frac{3\sqrt21-8}{25} = \frac{6\sqrt21-16}{5}$, so the ratio of $BP$ to $BC$ is $\frac{3\sqrt{21}+42}{100}$. This is the same as the ratio of $AE$ to $AD$, so $AD = 5\sqrt{21} \cdot \frac{100}{3\sqrt{21}+42} = \frac{20}{3} \cdot (2\sqrt{21}-3)$. 

 Dropping the altitude from $D$ to $AB$ and labelling the foot $Q$, $AQ$ can be (tediously) computed to be $\frac{20}{3} \cdot \sqrt{(2\sqrt{21}-3)^2-6^2} = \frac{20}{3} \cdot \sqrt{57-12\sqrt{21}} = \frac{20}{3} \cdot (6-\sqrt{21}) =\frac{120-20\sqrt{21}}{3}$, and subtracting this from $AB$ gives the length of $BQ$ equivalent to $CD$, which is our desired length, so our answer is $\frac{20\sqrt{21}-75}{3} \implies \boxed{\text{(A) }119}.$




\textbf{Solution 6}


                    \begin{figure}[H]
                        \centering
                        \includegraphics[width=0.5\textwidth]{Solution/AMC_10_10A_2025_20_Solution_Figure_4.png}
                    \end{figure}
                    
$F$ is the midpoint of $AD$, $DG \perp AB$. Notice that $\triangle OEF \sim \triangle DGA$, we have
 \[\frac{OE}{OF} = \frac{DG}{DA}\]Given the length of $CD$ is $g$, $OF$ is the midline, $CD = BG$, therefore $GA = 15 - g$, $DA = \sqrt{40^2 + (15 - g)^2}$, $OF = \dfrac{15 + g}{2}$
 \[\frac{OE}{OF} = \frac{DG}{GA}\]
 \[\cfrac{10}{\cfrac{15 + g}{2}} = \frac{40}{\sqrt{40^2 + (15 - g)^2}}\]
 \[2(15 + g) = \sqrt{40^2 + (15 - g)^2}\]
 \[3g^2 + 150g - 925 = 0\]
 \[g = \frac{-150 + \sqrt{22500 + 11100}}{6}\]
 \[g = \frac{20\sqrt{21}-75}{3}\]The answer is $20+21+75+3=\boxed{\textbf{(A)}~119}$




\textbf{Solution 7 (Pythag bash with same diagram as Solution 6)}

We start by making the right triangle as in Solution 6.$AG = (15 - g)$Since $OE$ is a radius to a tangent line, $\angle OEA = \angle OED = 90^{\circ}$. Drawing $OD$ and using the Pythagorean Theorem on triangle $OCD$ gives $OD = \sqrt{g^2 + 400}$. Using the Pythagorean Theorem on triangle $OED$ gives $ED = \sqrt{g^2 + 300}$.Using the Pythaogrean Theorem on triangle $OBA$ gives $OA = 25$Using the Pythagorean Theorem on triangle $OEA$ gives $EA = 5\sqrt{21}$.

 Since $BCDG$ is a rectangle, $BC = GD = 40$.

 Using the Pythagorean Theorem on triangle $AGD$ gives $(15-g)^2 + 40^2 = (5\sqrt{21} + \sqrt{g^2 + 300})^2$Expanding the squares, rearranging, and simplifying gives the quadratic $3g^2 + 150g - 925 = 0$.Completing the square and taking only the positive square root seems like the easiest method to find $g$ here
 \[3(g^2 + 50g) = 925\]
 \[3(g^2 + 50g + 625) = 925 + 3(625)\]
 \[(g+25) = \sqrt{\frac{2800}{3}} = \frac{20\sqrt{21}}{3}\]
 \[g = \frac{20\sqrt{21} - 75}{3}\]

 Finally, $a + b + c + d = 20 + 21 + 75 + 3 = \boxed{\textbf{(A) 119}}$



	\item (\textit{AMC 10 10A 2025}) A set of numbers is called \textit{sum-free} if whenever $x$ and $y$ are (not necessarily distinct) elements of the set, $x+y$ is not an element of the set. For example, $\{1,4,6\}$ and the empty set are sum-free, but $\{1,4,5\}$ is not. What is the greatest possible number of elements in a sum-free subset of $\{1,2,3,...,20\}$?


 $\textbf{(A) } 8 \qquad\textbf{(B) } 9 \qquad\textbf{(C) } 10 \qquad\textbf{(D) } 11 \qquad\textbf{(E) } 12$

 \textbf{Answer: }C



\textbf{Solution 1}

Let our subset be $\{11,12,13,...,20\}.$ If we add any element from the set $\{1,2,3,...,10\}$ to our current subset, we will have to remove at least one element from our subset. We cannot add 10 because if we do, the sum of 10 and itself is 20. Hence, the maximum size of our subset is $\boxed{\text{(C) }10}$.

 



\textbf{Solution 2 (Quicksolve)}

Notice that an odd number plus an odd number always results in an even number. Consider the subset of all odd numbers $S = \{1, 3, 5, ..., 19\}$. The addition of any even number will result in a violation of the rules, so the maximum number of elements in this subset is $\boxed{\text{(C) }10}$.

 \textbf{Note (Those who know)} This problem is essentially the same as \href{https://artofproblemsolving.com/wiki/index.php?title=2022\_AMC\_10B\_Problems/Problem\_14}{2022 AMC 10B Problem 14}.

 



\textbf{Solution 3}

Let $S = \{1, 2, 3, ..., 20\}$. We are looking for the largest possible sum-free subset $A \subseteq S$.

 First, we can show that a sum-free set of size 10 is possible.Consider the set $A_1 = \{11, 12, ..., 20\}$. The smallest possible sum of two elements (not necessarily distinct) from this set is $11 + 11 = 22$. Since $22 > 20$, no sum of two elements from $A_1$ can be in $A_1$. The size of $A_1$ is $20 - 11 + 1 = 10$.Thus, a sum-free set of size 10 does exist.

 Next, we prove that no sum-free set $A \subseteq S$ can have 11 or more elements.Let $m$ be the largest element in $A$.Since $A$ is sum-free and $m \in A$, for any $x \in A$, we must have $m-x \notin A$ (unless $x = m-x$, in which case $x \notin A$). This is because if both $x$ and $m-x$ were in $A$, their sum $x + (m-x) = m$ would be in $A$, which is a contradiction.

 We can partition the set $\{1, 2, ..., m\}$ into pairs that sum to $m$, plus any leftover elements.

 Case 1: $m$ is even. Let $m = 2k$.The numbers in $\{1, ..., m\}$ can be grouped as:$\{ (1, 2k-1), (2, 2k-2), ..., (k-1, k+1) \}$ along with the numbers $k$ and $m=2k$.

 
            \begin{itemize}
        \item $A$ contains $m$.
        \item $A$ cannot contain $k$, because $k+k = 2k = m$, which is in $A$.
        \item $A$ can contain at most one number from each of the $k-1$ pairs.
            \end{itemize}
            The maximum number of elements $A$ can contain from $\{1, ..., m\}$ is $1 + 0 + (k-1) = k = m/2$.Since $m \le 20$, the maximum size of $A$ is $m/2 \le 20/2 = 10$.

 Case 2: $m$ is odd. Let $m = 2k-1$.The numbers in $\{1, ..., m\}$ can be grouped as:$\{ (1, 2k-2), (2, 2k-3), ..., (k-1, k) \}$ along with $m=2k-1$.

 
            \begin{itemize}
        \item $A$ contains $m$.
        \item $A$ can contain at most one number from each of the $k-1$ pairs.
            \end{itemize}
            The maximum number of elements $A$ can contain from $\{1, ..., m\}$ is $1 + (k-1) = k$.Since $m = 2k-1$, we have $k = (m+1)/2$.The largest possible odd $m$ is 19. If $m=19$, $k = (19+1)/2 = 10$.So the maximum size of $A$ is $(m+1)/2 \le (19+1)/2 = \boxed{\textbf{(C) } 10}$.

 In all cases, a sum-free subset of $S$ can have at most 10 elements.Since we found a set with 10 elements, the greatest possible number of elements is $10$.

 



\textbf{Solution 4 (Another quicksolve)}

Note that selecting integers from $\{11, 12, ..., 20\}$ also satisfy the constraints of the problem. The minimum sum of any two members of the set is $11 + 11 = 22$ which is greater than any of the elements.

 



\textbf{Solution 5 (with quick proof)}

As all of the above solutions point out, the set $\{11, ..., 20\}$ satisfies the required condition, and so the greatest size of a sum-free set is no less than 10. How do we know that 11 isn't possible? Induction is our friend here. 

 We want to prove that the set $\{1, ..., 2n\}$ has a sum-free set of size no more than $n$. It is obvious for $n=1$. Suppose it is true for $n-1$. Now let $S$ be a sum-free subset of $\{1, ..., 2n\}$ and that it has at least $n+1$ elements. The intersection of $S$ and $\{1, ..., 2n-2\}$ is sum-free, so it has no more than $n-1$ elements. So we must include both $2n-1$ and $2n$ into $S$. If so, then $1$ and $n$ cannot belong to $S$. 

 We are now searching for $n-1$ more elements. But we can make $n-2$ 'pigeonholes': $\{2, 2n-2\}, \{3, 2n-3\}, ..., \{n-1, n+1\}$, seeing that they can provide no more than $n-2$ elements because each pigeonhole will have a sum of $2n$.

 


	\item (\textit{AMC 10 10A 2025}) A circle of radius $r$ is surrounded by three circles, whose radii are 1, 2, and 3, all externally tangent to the inner circle and externally tangent to each other, as shown in the diagram below.


 
                    \begin{figure}[H]
                        \centering
                        \includegraphics[width=0.5\textwidth]{Problem/AMC_10_10A_2025_22_Problem_Figure_1.png}
                    \end{figure}
                    


 What is $r$?


 $\textbf{(A) }\frac{1}{4}\qquad\textbf{(B) }\frac{6}{23}\qquad\textbf{(C) }\frac{3}{11}\qquad\textbf{(D) }\frac{5}{17}\qquad\textbf{(E) }\frac{3}{10}$

 \textbf{Answer: }B



\textbf{Solution 1}

\href{https://artofproblemsolving.com/wiki/index.php?title=Descartes\%27\_Circle\_Formula}{Descartes' Circle Formula} (curvatures $k_i = \frac{1}{r_i}$)
 \[k_4 = k_1 + k_2 + k_3 \pm 2\sqrt{k_1k_2 + k_2k_3 + k_3k_1}.\]

 For radii 1, 2, 3 we have
 \[k_1 = 1,\quad k_2 = \frac{1}{2},\quad k_3 = \frac{1}{3}.\]

 Compute the sum and the square-root term
 \[k_1+k_2+k_3 = \frac{11}{6},\qquad k_1k_2+k_2k_3+k_3k_1 = 1.\]

 Therefore
 \[k_4 = \frac{11}{6} \pm 2.\]

 Choose the plus sign for the small circle tangent externally to the three given circles
 \[k_4 = \frac{11}{6} + 2 = \frac{23}{6}, \qquad r_4 = \frac{1}{k_4} = \boxed{\textbf{(B) }\frac{6}{23}}.\]~Jonathanmodc

 



\textbf{Solution 2 (Descartes’ Theorem, Detailed Derivation)}

We are given the radii of three circles and are asked to find an external tangent to the radius of the circle.

 We can use Descartes’ Theorem to find it using curvatures (reciprocals of the radii). The curvatures are
 \[k_1 = 1, \quad k_2 = \frac{1}{2}, \quad k_3 = \frac{1}{3}.\]

 \href{https://artofproblemsolving.com/wiki/index.php?title=Descartes\%27\_Circle\_Formula}{Descartes' Circle Formula} states
 \[(k_1 + k_2 + k_3 + k_4)^2 = 2(k_1^2 + k_2^2 + k_3^2 + k_4^2).\]

 Plugging in, we get
 \[\left(1 + \frac{1}{2} + \frac{1}{3} + k_4\right)^2 = 2\left(1 + \frac{1}{4} + \frac{1}{9} + k_4^2\right).\]

 Simplifying,
 \[\left(\frac{11}{6} + k_4\right)^2 = 2\left(\frac{49}{36} + k_4^2\right).\]
 \[\frac{121}{36} + \frac{11}{3}k_4 + k_4^2 = \frac{49}{18} + 2k_4^2.\]
 \[k_4^2 - \frac{11}{3}k_4 - \frac{23}{36} = 0.\]

 Using the quadratic formula,
 \[k_4 = \frac{\frac{11}{3} \pm \sqrt{\frac{121}{9} - (-\frac{23}{9})}}{2} = \frac{\frac{11}{3} \pm \sqrt{\frac{144}{9}}}{2} = \frac{\frac{11}{3} \pm \frac{12}{3}}{2} = \frac{23}{6}, -\frac{1}{6}.\]

 Since the curvature is the reciprocal of the radius, the two tangent circles possible are
 \[r = \frac{1}{|k_4|} = \frac{6}{23} \text{ and } 6.\]

 The circle with radius 6 is the one that contains all of the three tangent circles inside it, so the answer is
 \[\boxed{(B)\ \frac{6}{23}}.\]

 Note: Instead of solving the quadratic directly, we can multiply it by $36$ and then use the rational root theorem to note that the solution must have a numerator of either $1$ or $23$, meaning that the denominator of the radius must be $1$ or $23$. It can't be $1$, since then the radius would be greater than $1$ (which is impossible, looking at answer choices), so the denominator is $23$, and there is only one answer choice whose denominator is $23$, $\boxed{(B)\ \frac{6}{23}}$




\textbf{Solution 3(Heron's)}

Let the center of the circle radius 2 be $A$, radius 1 be $B$, radius 3 be $C$. Let the unknown circle's center be $P$ and radius be $r$.

 We know $\overline{AB}=2+1=3, \overline{AC}=3+2=5, \overline{BC}=1+3=4.$

 Thus, $\triangle ABC$ is a 3-4-5 triangle, where $\angle B=90^\circ$.

 Let $X,Y$ be on $\overline{AB},\overline{BC}$ respectively such that $\angle{PXB}=\angle{PYB}=90^\circ$.

 By Heron's formula, the area of $\triangle BPC$ is 
 \[\sqrt{s(s-\overline{BP})(s-\overline{CP})(s-\overline{BC})}=\sqrt{(4+r)\cdot 1\cdot 3\cdot r}=\sqrt{3r^2+12r}.\]By using the simpler area formula, we can find $\overline{PY}$.
 \[\overline{PY}=\frac{[BPC]}{\overline{BC}/2}=\frac{\sqrt{3r^2+12r}}{2}.\] 

 Since $\angle{PXB}=\angle{B}=\angle{PYB}=90^\circ$, $XPYB$ is a rectangle, and we can say 
 \[\overline{PX}^2=\overline{BP}^2-\overline{PY}^2,\]or
 \[\overline{PX}^2=(1+r)^2-\frac{3r^2+12r}{4}=\frac{r^2-4r+4}{4}\]

 \[\Rightarrow \overline{PX}=\frac{|r-2|}{2}.\]

 We then find the area of $ABP$ via Heron's formula, and again, by the simple area formula we have that
 \[\overline{PX}=\frac{[ABP]}{\overline{AB}/2}=\frac{2\sqrt{2r^2+6r}}{3}\]

 Now we can equate the two and solve the equation:

 
 \begin{align*} \frac{|r-2|}{2}&=\frac{2\sqrt{2r^2+6r}}{3} \\ \Rightarrow 3(|r-2|)&=4\sqrt{2r^2+6r} \\ \Rightarrow 9r^2-36r+36&=32r^2+96r \\ \Rightarrow 23r^2+132r-36&=0 \end{align*}
Utilizing the quadratic formula, the only positive root is $r=\frac{6}{23}=\boxed{B}.$

 Note: By the rational root theorem, a rational root of the final quadratic must have denominator $1$ or $23$ in simplest form, which eliminates all choices except B.

 Edited by GarudS

 



\textbf{Solution 4 (Variation of Solution 3)}

Follow the same starting steps as solution 3 (drawing right triangle ABC connecting the centers of the three larger circles and connecting points A, B, and C to the center of the smaller circle). Using Heron's formula, we can easily find the areas of the three smaller triangles that make up triangle ABC in terms of $r$. Since the sum of these terms is an integer, or 6, we can theorize that the square roots have to simplify out into rational numbers. Out of all the answer choices, $\boxed{(B)}$, or $\dfrac{6}{23}$, is the only one that works.(will add more later)

 



\textbf{Solution 5 (Coordinate Geometry)}

Let the center of the circle radius 2 be $A$, radius 1 be $B$, radius 3 be $C$. Let the unknown circle's center be $P$.Since $AB^2+BC^2=AC^2$, $\triangle ABC$ is a right triangle, with right angle at $B$.Let $B=(0, 0), A=(3, 0), C=(0, 4),$ and $P=(x, y)$. Since the large circles are tangent to the small one, we have, if the radius of the smallest circle is $r$,
 \[x^2+y^2 = (r+1)^2 = r^2+2r+1 \\\]
 \[(x-3)^2+y^2 = (r+2)^2 = r^2+4r+4 \\\]
 \[x^2+(y-4)^2 = (r+3)^2 = r^2+6r+9\]Subtracting (1) from (2) and (1) from (3), we get
 \[-6x+9=2r+3 \implies -3x=r-3\\\]
 \[-8y+16=4r+8 \implies -2y=r-2\]To avoid fractions, we multiply (1) to get $36r^2+72r+36=36x^2+36y^2=4(r-3)^2+9(r-2)^2 = 13r^2-60r+72$. Combining like terms, we have $23r^2+132r-36$. Solving via quadratic formula, we get $\boxed{\frac{6}{23}}$.

 Note: Like the last solution, you can just notice that $B$ is the only answer choice with 23 in the denominator.

 \textbf{A note about factoring $23r^2+132r-36$} You don't need to save a lot of time by abusing multiple choice to cheese the root of $23r^2+132r-36$.

 $23r^2+132r-36$ = $(Ax+B)(Cx+D)$. Since $23$ is prime, $B=1$ and $D=23$, WLOG.

 Theorem you can prove and learn: in $Px^2 + Qx + R =  (Ax+B)(Cx+D)$: if $n \vert Q$ and $n^2 \vert R$, then $n\vert C$ and $n \vert D$.(This is equivalent to substituting $y=x/n$.)Symmetrically, if $n \vert Q$ and $n^2 \vert P$, then $n\vert A$ and $n \vert B$.(This is equivalent to substituting $y=nx$.)

 Back to our problem: Since $6 \vert 132$  and $6^2 \vert 36$, then $6 \vert A$ and $6 \vert B$. $6 \cdot 6 =36$ so the factoring is done, up to placing the minus sign: $(23r-6)(r+6)$.

 Considering the signed-distance geometry of the original equations in $r$, the positive root  $\boxed{6/23}$ is the internally tangent circle's radius, and the absolute value $6$ of negative root $-6$ is the externally tangent circle's radius.

 



\textbf{Solution 6 (Trig, LoC)}

Start with the same setup as Solution 3. Now, set $\angle PBC$ as $\theta$. Law of Cosine on $\triangle PBC$, you get:

 
 \[\cos(\theta )=\frac{PC^2-PB^2-BC^2}{2\cdot PB\cdot BC}\]

 \[=\frac{(3+r)^2-(1+r)^2-16}{8+8r}\]

 \[=\frac{r-2}{2+2r}\]

 Similarly, use LoC on $\triangle APB$:

 
 \[\cos(90^\circ -\theta )=\sin(\theta )=\frac{PA^2-PB^2-AB^2}{2\cdot PB\cdot AB}\]

 \[=\frac{(2+r)^2-(1+r)^2-9}{6+6r}\]

 \[=\frac{r-3}{3+3r}\]

 Now, Pythagorean Identity,

 
 \[\sin ^2(\theta)+\cos ^2(\theta)=1\]

 \[\Rightarrow \left( \frac{r-2}{2+2r}\right) ^2+\left( \frac{r-3}{3+3r}\right) ^2=1\]

 \[\Rightarrow \frac{13r^2-60r+72}{36x^2+72x+36}=1\]

 \[\Rightarrow 23r^2+132x-36=0\]

 Solving, you get $r=\boxed{\frac{6}{23}}$.

 



\textbf{Solution 7 (LAST RESORT)}

DISCLAIMER: This solution is extremely risky and should only be used if you want to get this problem correct but you either have no time or you can't solve it. Also, if this were the case, it would be a better idea to leave the answer blank and get a guaranteed 1.5 points.

 
Rulers are permitted on the AMC 10. So, you can measure the diameter of the circle with radius $1$. In this case, the circle has an approximate diameter of $2.4$ centimeters. Thus, it's radius is approximately $1.2$ centimeters. Now we can measure the diameter of the smallest circle, which is approximately $0.6$ centimeters (a little more, keep that in mind), so it's radius is approximately $0.3$ centimeters (The real value is a little more than that, just keep that in mind). If $1.2$ centimeters correspond to 1 unit, then $0.3$ centimeters has to correspond to $\frac{1}{4}$th of a unit. So we check which answer choice is closest to $\frac{1}{4}$, BUT not exactly $\frac{1}{4}$, since we noted that the measurement of the radius of the smallest circle is not exactly $0.3$ centimeters. So, from this, we find that $\boxed{\textbf{(B) }\frac{6}{23}}$ has the smallest error compared to $\frac{1}{4}$.

 
Note: Although MAA claims that the diagrams are not drawn to scale, they almost certainly are (there are a few exceptions but this is the general, most likely case).

 



\textbf{Solution 8 (Another version of Kissing Circles Theorem)}

Yes, I am aware this is the same as solution 2. Now, if you ask me why this is here, before I made this I thought solution 2 used the same version of the Kissing Circles Theorem like solution 1. I wouldn't remove this as it is faster than solution 2 in an essence, but to be fair I rlly don't know why I made this :sob:).

 Another relation between four curvatures via Kissing Circles Theorem is

 
 \[(k_1 + k_2 + k_3 + k_4)^2 = 2(k_1^2 + k_2^2 +k_3^2 + k_4^2)\]

 Then $k_1 = 1$, $k_2 = \frac{1}{2}$, $k_3 = \frac{1}{3}$. Find $k_4$.

 Substitution gives

 
 \[(\frac{11}{6} + k_4)^2 = 2(1+\frac{1}{4} + \frac{1}{9} + k_4^2)\]
 \[\frac{121}{36} + \frac{11}{3}k_4 + k_4^2 = \frac{49}{18} + 2k_4^2\]
 \[k_4^2 - \frac{11}{3}k_4 -\frac{23}{36} = 0\]
 \[36k_4^2 -132k_4 - 23 = 0\]

 The two values $k_4$ are $k_4 = \frac{23}{6}$ and $k_4 = -\frac{1}{6}$. Negative curvatures go beyond Plane Euclidean Geometry's definition of external tangencies. Thus $k_4 = \frac{23}{6}$. It follows that $r_4 = \boxed{\frac{6}{23}}$.

 



\textbf{Solution 9 (More in depth of Solution 4)}

Let $r$ be the radius of the center circle. We wish to find the value of $r$.First, connect the centers of the circles with radius 1 and radius 2, radius 1 and radius 3, and radius 2 and radius 3. This forms a triangle with sides 3, 4, and 5. Note that $3^2 + 4^2 = 5^2$. By the converse of the Pythagorean theorem, this is a right triangle. So the area of this triangle is $\frac{(3)(4)}{2} = 6$. Now draw the radii from the center circle connecting to all of the other 3 circle centers. This forms 3 triangles, which we can find the area of in terms of $r$ with Heron's Formula.

 Triangle 1: This triangle has sides $r+1, r+2, 3$. The area of this triangle is $\sqrt{(r+3)(2)(1)(r)}$.

 Triangle 2: This triangle has sides $r+1, r+3, 4$. The area of this triangle is $\sqrt{(r+4)(3)(1)(r)}$

 Triangle 3: This triangle has sides $r+2, r+3, 5$. The area of this triangle is $\sqrt{(r+5)(3)(2)(r)}$

 Since these three triangles make up the area of the whole triangles, we have the equation $\sqrt{(r+3)(2)(1)(r)} + \sqrt{(r+4)(3)(1)(r)} + \sqrt{(r+5)(3)(2)(r)} = 6$. Now instead of solving this disgusting equation, we can just plug in the answer choices given. Checking all of them gives $r = \boxed{\textbf{(B)}\frac{6}{23}}$, which is our final answer.

 


 


	\item (\textit{AMC 10 10A 2025}) Triangle $\triangle ABC$ has side lengths $AB = 80$, $BC = 45$, and $AC = 75$. The bisector of $\angle B$ and the altitude to side $\overline{AB}$ intersect at point $P.$ What is $BP$?


 $\textbf{(A)}~18\qquad\textbf{(B)}~19\qquad\textbf{(C)}~20\qquad\textbf{(D)}~21\qquad\textbf{(E)}~22$

 \textbf{Answer: }D



\textbf{Solution 1}

Let $CD \perp AB$ with foot $D$. Right triangles $ACD$ and $BCD$ give $AC^2 = AD^2+CD^2$, $BC^2 = BD^2+CD^2$, $AC^2-BC^2 = AD^2-BD^2 =(AD-BD)(AD+BD)$. 

 Since $AD+BD = AB = 80$ and $AC^2-BC^2 = 75^2-45^2 = 3600$, we get the equation $3600 = 80(AD-BD)$. This equation simplifies to $45 = AD - BD$. We can solve the system of equations $AD + BD = 80$ and $AD - BD = 45$ easily via elimination, and we get $AD = \frac{125}{2}$, $BD = \frac{35}{2}$. $CD^2 = AC^2-AD^2 = 75^2-\left(\frac{125}{2}\right)^2 = \frac{6875}{4}$, $CD = \frac{25\sqrt{11}}{2}$. 

 By Angle Bisector Theorem, $\frac{DP}{PC} = \frac{DB}{BC} = \frac{\frac{35}{2}}{45} = \frac{7}{18}$, $PC = CD-DP$ thus, $18DP = 7(CD-DP)$, $25DP = 7CD$, $DP = \left(\frac{7}{25}\right)CD = \left(\frac{7}{25}\right)\left(\frac{25\sqrt{11}}{2}\right) = \frac{7\sqrt{11}}{2}$. $BP^2 = BD^2+DP^2 = \left(\frac{35}{2}\right)^2+\left(\frac{7\sqrt{11}}{2}\right)^2 = \frac{1225}{4}+\frac{49(11)}{4} = \frac{1764}{4} = 441$, thus $BP = \boxed{\text{(D) }21}.$




\textbf{Solution 2 (Law of Cosines)}

Scale this down to a $9-15-16$ triangle (we will multiply the result by $5$ in the end).Note that $9^2+16^2-18\cdot16\cdot\cos(\angle B) = 15^2$, so $112 = 18 \cdot 16 \cdot \cos(\angle B)$, which simplifies to $\frac{7}{18} = \cos(\angle B)$. Then $\cos(\frac{\angle B}{2}) = \sqrt{\frac{1+\frac{7}{18}}{2}} = \sqrt{\frac{25}{36}} = \frac{5}{6}$ (positive root since the angle is acute). Therefore, we have $\frac{BD}{BP} = \frac{5}{6}$, assuming that $D$ is the foot of the altitude.There are many ways to proceed from here to find $BD$. Note that by Heron's formula, the area of the scaled-down triangle is $\sqrt{(20)(4)(5)(11)} = 20\sqrt{11}$. Therefore, $CD = \frac{5}{2}\sqrt{11}$. Using Pythagorean Theorem, we get $BD= \sqrt{81-\frac{25 \cdot 11}{4}} = \sqrt{\frac{324-275}{4}} = \frac{7}{2}$. Therefore, we get $\frac{\frac{7}{2}}{BP} = \frac{5}{6}$, so $BP = \frac{21}{5}$, and we scale up by $5$ to get $\boxed{\text{(D) }21}$.~ScoutViolet

 



\textbf{Solution 3 (Stewarts)}

Let the foot of the altitude coming from $C$ on segment $AB$ be $D$. Using the fact that $CD$ is a common leg in right triangles $\triangle CDB$ and $\triangle CDA$, we have 
 \[45^2 - BD^2 = 75^2 - (80-BD)^2.\] Expanding gives 
 \[45^2=75^2-80^2+160BD,\] so $BD=\frac{35}{2}.$ Let the foot of the angle bisector from $B$ to $AC$ be point $E$. Since $BE$ is the angle bisector of $\angle CBA$, we can use the angle bisector theorem. This gives 
 \[\frac{45}{CE}=\frac{80}{75-CE},\] so $CE=27$ and $AE=75-27=48$. Now we can use Stewart’s Theorem to find $\overline{BE}$. We have 
 \[(48)(45)^2+(27)(80)^2=(27)(48)(75)+75BE^2.\] To simplify this expression, just divide by the greatest common divisor and solve from there. In the end, we get $BE=48$. Let $BP=x$, so $PE=48-x$. Draw the altitude from $E$ down to $AB$. Let the foot of this altitude be $F$. Since $EF || CD$, we have $\triangle AFE \sim \triangle ADC$. Hence, we can write the equation 
 \[\frac{48}{75}=\frac{AF}{\frac{125}{2}}.\] Solving gives $AF=40$, so $FD=\frac{125}{2}-40=\frac{45}{2}$. Since $PD || EF$, we also have $\triangle BDP \sim \triangle BFE$, so we have 
 \[\frac{x}{48}=\frac{\frac{35}{2}}{40}.\] Solving for $x$ gives $\boxed{x=21}$ or $\boxed{\text{D}}.$




\textbf{Solution 4 (Rulerbash)}

Preface: When we have a problem as such, involving a simple diagram with minimal instructions, I use a method I named "rulerbash". Rulerbash should only be used in specific cases and as a last resort, mainly in the event of a time crunch or a difficult problem. 

 Start with the longest side, drawing a line with a length of $8\,\text{cm}$ ($AB$). Then, using a compass, draw 2 circles centered around points $A$ and $B$, $7.5\,\text{cm}$ and $4.5\,\text{cm}$ radii respectfully. At the point of intersection of these 2 circles, we have point $C$, completing a perfectly scaled drawing of $\triangle ABC$. (Note the circles are not necessary with a bit of trial and error with the side lengths, they simply offer a way to get it done first try). 

 
                    \begin{figure}[H]
                        \centering
                        \includegraphics[width=0.5\textwidth]{Solution/AMC_10_10A_2025_23_Solution_Figure_1.png}
                    \end{figure}
                    

 


 
After this, dropping the altitude to $AB$ is simple with a ruler and careful placement, and angle bisector can be estimated quite accurately. (Note: Alternately, a compass can be used to accurately place the altitude and the bisector, which are both trivial constructions ~ZingiberiMaracandae)

 
                    \begin{figure}[H]
                        \centering
                        \includegraphics[width=0.5\textwidth]{Solution/AMC_10_10A_2025_23_Solution_Figure_2.png}
                    \end{figure}
                    

 After all of this, we can reuse our ruler and measure $BP = 2.1\,\text{cm}$, and using our scale of $80=8\,\text{cm}$, our final answer is $\boxed{\text{(D) }21}.$

 Other note: you kind of have to draw the angle bisector really accurately for this to work (I tried this)~ Sakura\_kitty

 



\textbf{Solution 5 (No Trig)}

Let \(D\) be the intersection of the altitude from \(C\) to \(AB\). To simplify calculations, divide all side lengths by \(5\), and multiply by \(5\) again at the end.  First, we use Heron’s Formula, \(\sqrt{s(s-a)(s-b)(s-c)}\), to find the area. Let \([ABC]\) denote the area of \(\triangle ABC\). By Heron’s Formula,  
 \[[ABC] = \sqrt{20 \cdot 5 \cdot 4 \cdot 11} = 20\sqrt{11}.\]  Next, we find the altitude \(CD\) using the formula for the area of a triangle, \(\frac{1}{2}bh = \text{area}\):  
 \[\frac{1}{2}(16)(CD) = 20\sqrt{11} \quad \Rightarrow \quad CD = \frac{5\sqrt{11}}{2}.\]  We can use the Pythagorean Theorem in \(\triangle CDB\) to find \(DB\):  
 \[DB^2 + \frac{25 \cdot 11}{4} = 81 \quad \Rightarrow \quad 4DB^2 + 275 = 324 \quad \Rightarrow \quad DB^2 = \frac{49}{4},\]so \(DB = \frac{7}{2}\).  Next, we use the Angle Bisector Theorem to find \(PD\). Let \(x = PC\) and \(y = PD\). Since \(x + y = \frac{5\sqrt{11}}{2}\), we have \(x = \frac{5\sqrt{11}}{2} - y\).  From the given ratio,  
 \[\frac{9}{x} = \frac{7}{2y} \quad \Rightarrow \quad 18y = 7x.\]  Substituting \(x = \frac{5\sqrt{11}}{2} - y\),  
 \[18y = 7\left(\frac{5\sqrt{11}}{2} - y\right) \quad \Rightarrow \quad 25y = \frac{35\sqrt{11}}{2},\]so \(y = \frac{7\sqrt{11}}{10}\).  

 Now, using the Pythagorean Theorem again to find \(BP\):  
 \[\frac{49 \cdot 11}{100} + \frac{49}{4} = BP^2 \quad \Rightarrow \quad 100BP^2 = 49(11 + 25) = 49 \cdot 36,\]so \(BP = \frac{42}{10} = \frac{21}{5}.\)  

 Finally, multiplying the side lengths by \(5\) again gives \(BP = 21.\), or $\boxed{\text{D}}.$




\textbf{Solution 6 (Similar Triangle)}


                    \begin{figure}[H]
                        \centering
                        \includegraphics[width=0.5\textwidth]{Solution/AMC_10_10A_2025_23_Solution_Figure_3.png}
                    \end{figure}
                    

 Due to Angle Bisector Theorem
 \[\frac{CE}{AE} = \frac{BC}{BA} = \frac{45}{80} = \frac{9}{16}\]
 \[CE = AC \times \frac{9}{25} = 27\]
 \[AE = AC \times \frac{16}{25} = 48\]Notice that $\dfrac{CE}{CB} = \dfrac{CB}{CA} = \dfrac{3}{5}$, $\triangle CBE \sim \triangle CAB$, $\angle A = \angle CBE = \angle EBF$, so $\triangle ABE$ is isosceles triangle, hence
 \[AF = BF\]Notice that $AF:FD = AE:EC = 16:9$, $DB = FB - FD$, then 
 \[AF:DB = 16:7\]

 Notice that $\triangle PBD \sim \triangle EAF$, $PB = AE \times \dfrac{BD}{AF} = 48 \times \dfrac{7}{16} = \boxed{\text{(D) }21}.$




\textbf{Solution 7 (Pythagoras+angle bisector theorem)}

Let $CD$ be the altitude to side $AB$, and let $E$ be the point of intersection of the angle bisector of $\angle B$. Let $P$ be the point of intersection of the angle bisector and the altitude.

 By Pythagoras, $CD^2=BC^2-BD^2=AC^2-AD^2$

 We let $BD = x$, so $AD=80-x$

 Now we have

 
 \[45^2-x^2=75^2-(80-x)^2\]

 (This looks complicated, but we can see that we are going to have two $-x^2$ terms on both sides, so it cancels out)

 Simplifying and rearranging gives us $80-2x=45$, so $x=\frac{35}{2}$

 Therefore, $BD=\frac{35}{2}$, and $AD=\frac{125}{2}$. This gives us $CD=\frac{25}{2}\sqrt{11}$

 
Now, note that in triangle $BDC$, $BP$ bisects $\angle DBC$

 By the angle bisection theorem, we have$\frac{BD}{DP}=\frac{BC}{CP}$

 which gives us $\frac{DP}{CP}=\frac{BD}{BC}=\frac{7}{18}$

 We have $DP=\frac{7}{25}*CD=\frac{7}{2}\sqrt{11}$ and $CP=\frac{18}{25}*CD=\frac{18}{2}\sqrt{11}$

 So

 $BP^2=BD*BC-DP*CP=\frac{35}{2}*45-\frac{7}{2}*\frac{18}{2}*11=441$

 Therefore, $BP=21$

 (This is my first solution, pls point out any mistakes in math or LATEX). ~backtosq-1

 \textbf{Alternative Ending to Solution 7} As an alternative ending equation to the solution above, I used the Pythagorean Theorem again, giving us: $BP^2=BD^2+PD^2=(\dfrac{35}{2})^2+(\dfrac{7\sqrt{11}}{2})^2=(\dfrac{1225}{4})+(\dfrac{539}{4})=\dfrac{1764}{4}=441$

 Thus, $BP=\boxed{\textbf{(D)}~21}$




\textbf{Solution 8 (Pythagora's + Double Angle Identity)}

Since all the lengths share a factor of $5$, divide by $5$ and multiply it back at the end for smaller numbers. Drop the altitude from vertex $C$ to point $D$ on $AB$. Then comparing the altitude by Pythagora's on both sides gives 
 \[9^2-x^2=15^2-(16-x)^2=15^2-16^2+32x-x^2=-31+32x-x^2\]Cancelling the $-x^2$ then solving gives us $x=\frac{112}{32}=\frac{7}{2}$. Since $\triangle BDC$ is a right triangle with $\angle B=2\theta$, we find $\cos(2\theta)=\frac{\frac{7}{2}}{9}=\frac{7}{18}$.

 Let $\theta=\angle PBA=\angle PBC$. Because $\angle B$ is opposite to the leg $9<15$ and not the hypotenuse, we must have that $2\theta$ is acute, likewise for $\theta$. Hence, $\cos \theta>0$.

 Thus, using cosine definition on $\triangle BDP$ and the double angle identity $\cos 2\theta=2\cos^2 \theta-1$, we get
 \[\cos \theta=\frac{7/2}{BP}=\sqrt{ \frac{1+\cos 2\theta}{2} }=\sqrt{ \frac{1+\frac{7}{18}}{2} }=\sqrt{ \frac{25}{36} }=\frac{5}{6}\]Finally, solving for $BP=\frac{7}{2\cos\theta}=\frac{7}{2}\cdot \frac{6}{5}=\frac{21}{5}$. Scaling back to the original diagram by $5$ gives $\boxed{\textbf{(D)}~21}$.

 



\textbf{Solution 9(Law of Cosines+Angle Bisector Length Theorem)}

Use Law of Cosines in $\triangle ABC$:
 \[\cos B=\frac{80^2+45^2-75^2}{2\cdot 80\cdot 45}=\frac{7}{18}\]Call $E$ the intersection of the angle bisector of $B$ and line segment $AC$, and let $CD \perp AB$, $D=CD\cap AB$, therefore $P=CD\cap BE$. In right triangle $\triangle BCD$, use the definition of cosine:
 \[\cos B=\frac{BD}{BC}=\frac{7}{18},\sin B=\sqrt{1-\cos^{2}B}=\frac{5\sqrt{11}}{18}\]Therefore 
 \[BD=\frac{35}{2}, CD=\frac{25\sqrt{11}}{2}\]Use Angle Bisector Theorem in $\triangle BCD$: 
 \[\frac{BD}{BC}=\frac{PD}{PC}=\frac{7}{18}\]Therefore, 
 \[PD=\frac{7}{25}CF=\frac{7\sqrt{11}}{2}, PC=\frac{18}{25}CF=9\sqrt{11}\]Use Angle Bisector length formula:
 \[BP=\sqrt{BD\cdot BC-PD\cdot PC} =\boxed{\textbf{(D)} \ 21}\]




\textbf{Solution 10}

Let the point where the altitude to $AB$ meets $AB$ be $D$. We can compute $D$ easily. Let $BD = a \implies AD = 80 - a$. So

 $45^2 - a^2 = 75^2 - (80 - a)^2 \implies 45^2 - a^2 = 75^2 - 80^2 + 160a - a^2 \implies 45^2 - 75^2 + 80^2 = 160a \implies -30 \cdot 120 + 80 \cdot 80 = 10 \cdot 16a \implies -360 + 640 = 16a \implies a = \frac{280}{16} = \frac{35}{2}$. 

 Now let $\angle DBP = \theta$. Note this is an angle bisector so we can find $\cos(\angle ABC)$ where $\theta = \frac{1}{2} \angle ABC$. So we use Law of Cosines to get

 $80^2 + 45^2 - 160 \cdot 45 \cos(2\theta) = 75^2 \implies 16^2 + 9^2 - 32 \cdot 9 \cdot \cos(2\theta) = 15^2 \implies \cos(2\theta) = \frac{112}{32 \cdot 9} = \frac{7}{18}$. Now note

 $\cos(2\theta) = 2\cos^2(\theta) - 1 = \frac{7}{18} \implies \cos^2(\theta) = \frac{25}{36} \implies \cos(\theta) = \frac{5}{6}$. 

 Finally, by $SOH-CAH-TOA$ note that $\cos(\theta) = \frac{BD}{BP} = \frac{5}{6}$ and we know $BD = \frac{35}{2}$ so 

 $BP = \frac{6}{5} \cdot \frac{35}{2} = \boxed{21}$.

 


	\item (\textit{AMC 10 10A 2025}) Call a positive integer \textit{fair} if no digit is used more than once, it has no 0s, and no digit is adjacent to two greater digits. For example, $196$, $23$, and $12463$ are fair, but $1546$, $320$, and $34321$ are not fair. How many fair positive integers are there?
$\textbf{(A) } 511 \qquad\textbf{(B) } 2584 \qquad\textbf{(C) } 9841 \qquad\textbf{(D) } 17711 \qquad\textbf{(E) } 19682$

 \textbf{Answer: }C



\textbf{Solution 1}

To satisfy the conditions, a fair integer must have no digit be a local minimum. Let's say we have $n$ distinct digits, with each digit being a number from $1$ to $9$. To create a fair integer, we begin by placing the largest digit. For the second-largest digit, we can either place this digit to the right or to the left of the string already created. We have these $2$ options for the third-largest digit, and so on. Therefore, there are $2^{n-1}$ valid permutations to create a fair integer.

 We must also choose which digits will be in the permutation. If you are creating an $n$-digit long fair integer, there are $\binom{9}{n}$ ways to pick which digits will be in the number.

 Therefore, for each $n \in \{1,2,\dots, 9\}$, the number of fair integers of length $n$ is:
 \[\binom{9}{n} \cdot 2^{n-1}.\]Summing over all $n$:
 \[\sum_{n=1}^9{\binom{9}{n} \cdot 2^{n-1}}=\frac{1}{2}\left(\sum_{n=0}^9{\binom{9}{n}}2^n -1\right)=\frac{1}{2}\left((1+2)^9 -1 \right) = \frac{1}{2}(19682) = \boxed{9841}.\]Note that the Binomial Theorem was used to equate
 \[\sum_{n=0}^9{\binom{9}{n}}2^n = (1+2)^9.\]~lprado

 



\textbf{Solution 2}

Note every fair number will have an increasing string of digits, a maximum digit, then a decreasing string of digits. This is because if it decreases then increases, then the digit in the middle will be less than its adjacent digits.

 Let $n$ be the maximum digit. For each number $i<n$, we may either place $i$ before $n$, after $n$, or choose not to include it. Note this process will result in a unique number for every case, as the numbers before $n$ must be in increasing order, and the numbers after $n$ must be in decreasing order. Therefore, for each number $n$, we have $3^{n-1}$ cases.

 Since $n \in \{1,2,\cdots9\}$, we have:
 \[\sum_{n=1}^9 3^{n-1}=\frac{3^9-1}{3-1}=\boxed{9841}.\]




\textbf{Solution 2.5 (Cheese, Goldilocks Zone)}

If you consider options D and E too big, and A and B too small, then C is your answer.

 $\boxed {\text {(C) } 9841}$Also E is twice C so C is your answer 

 - Ameen Patel IrvingtonHighSchoolMathGeeks-Edited by TFEA

 



\textbf{Solution 3 (Recursion)}

Let $f(n)$ be the number of fair integers given an arbitrary set of $n$ digits. Let $a$ be the smallest of these digits. Notice that $a$ is either the first or last digit, as if it were any other digit, the two digits surrounding it would both be greater. Then, notice that if the remaining $n-1$ slots form a fair number, so does the number when $a$ is appended. Therefore, $f(n) = 2f(n-1)$. From here, we may proceed with the calculation in Solution $1$ to get the answer of $\boxed {\text {(C) } 9841}$




\textbf{Solution 4}

For any fair integer $n$, we write some $9$-digit base-$3$ number(s) $f(n)=\left(\overline{d_1 d_2 \ldots d_9}\right)_3$ from left to right as follows:

 
            \begin{itemize}
        \item $d_k=0$ if $k$ is not in $n$;
        \item $d_k=1$ if $k$ is the to the left of all digits greater than $k$ in $n$;
        \item $d_k=2$ if $k$ is the to the right of all digits greater than $k$ in $n$.
            \end{itemize}
            Note that the fairness ensured one and only one of the above three cases happens. 

 Now if $k$ is the largest digit in $n$ then $d_k$ can be either $1$ or $2$.

 An example: If $n=136852$, then $f(n)=\overline{121021010}_3$ or $\overline{121021020}_3$. 

 It's thus easy to see that there is a $1$-to-$2$ mapping from all fair integers to all $9$-digit base-$3$ nonzero numbers. Therefore the answer is $\frac{3^9-1}{2}$.

 


	\item (\textit{AMC 10 10A 2025}) A point $P$ is chosen at random inside square $ABCD$. The probability that $\overline{AP}$ is neither the shortest nor the longest side of $\triangle APB$ can be written as $\frac{a + b \pi - c \sqrt{d}}{e}$, where $a, b, c, d,$ and $e$ are positive integers, $\text{gcd}(a, b, c, e) = 1$, and $d$ is not divisible by the square of a prime. What is $a+b+c+d+e$?


 $\textbf{(A) }25 \qquad \textbf{(B) }26 \qquad \textbf{(C) }27 \qquad \textbf{(D) }28 \qquad \textbf{(E) }  29 \qquad$

 \textbf{Answer: }A



\textbf{Solution 1 (Geometry)}


                    \begin{figure}[H]
                        \centering
                        \includegraphics[width=0.5\textwidth]{Solution/AMC_10_10A_2025_25_Solution_Figure_1.png}
                    \end{figure}
                    

 Say WLOG

 


\end{enumerate}